\chapter{Supported compositions and outer collision classes}
\label{ch:vol2-supported-compositions}

Three quadratic current fields have a connected Gaussian correlation
with a pole on each pair diagonal. Its triangle of contractions already
lies beyond a forest calculation. A second operator product must act on
these poles and on their residue attachments, as well as on the states
left by the first product. The coefficient modules of
Chapter~\ref{ch:vol2-coefficients} provide the domains for those maps.

We begin with the finite Wick formula and its triangle current. Mixed
Green kernels give a second construction, in which each edge has a point
source and a homotopy must retain contact terms. Ordered collision
complexes then show which classes disappear when only consecutive
cuts are kept. The final comparison concerns compositions on supported
field coefficients, with the lower homotopies fixed.

Unless another base is specified, $k$ is a field of characteristic
zero. Polynomial states contain finitely many jets; infinite Taylor
series occur only in the indicated formal variables. Analytic currents
are taken over $\mathbb C$. Complexes are cohomological, with
$K[1]^n=K^{n+1}$, and the differential of a homogeneous map is
$\delta f=d_{\mathrm{out}}f-(-1)^{|f|}fd_{\mathrm{in}}$.

\section{Finite fields and the Gaussian triangle}
\label{sec:vol2-wick-fields}

Allow the Heisenberg level to be the indeterminate $\lambda$. The
polynomial representation of \eqref{eq:vol2-fock-currents} is then
defined over $k[\lambda]$, including at $\lambda=0$:
\[
 V=k[\lambda][x_1,x_2,\ldots],\qquad
 J_{-n}=L_{x_n},\quad J_n=\lambda n\partial_{x_n}\quad(n>0),
 \qquad J_0=0.
\]
Set $J=x_1$, $T=\sum_{n\geq1}n x_{n+1}\partial_{x_n}$ and
$j_r=T^rJ=r!x_{r+1}$. These are raw, undivided jets.
Write $J(z)=\sum_{m\in\mathbb Z}J_mz^{-m-1}$.
For a polynomial $F(j_0,j_1,\ldots)$, define its field $Y(F,z)$
by replacing $j_r$ with $\partial_z^rJ(z)$ and putting every
creation operator $J_{-n}$ to the left of every annihilation
operator $J_n$. The vacuum field is $Y(1,z)=1$.

The commutator of an annihilation part with a creation part is
the expansion of $\lambda/(z-w)^2$ in nonnegative powers of $w/z$.
Consequently the contraction between raw jets is
\begin{equation}\label{eq:vol2-jet-contraction}
 K_{rs}(z,w)
 =\partial_z^r\partial_w^s\frac{\lambda}{(z-w)^2}
 =\frac{\lambda(-1)^r(r+s+1)!}{(z-w)^{r+s+2}}.
\end{equation}
It is a meromorphic coefficient. A chosen ordering of the formal
variables specifies its Laurent expansion.

For polynomial inputs $F_i$ at distinct labels $i=1,\ldots,n$, let
$j_r^{(i)}$ be independent copies of the jet variables and put
\[
 \Gamma_{ij}
 =\sum_{r,s\geq0}K_{rs}(z_i,z_j)
       \partial_{j_r^{(i)}}\partial_{j_s^{(j)}}\qquad(i<j).
\]
After these operators act, colons mean that the surviving variables
are replaced by their fields at the stated positions and all resulting
creation operators are moved to the left, without further contractions.

\begin{proposition}[Polynomial Wick evaluator]
\label{prop:vol2-wick-evaluator}
On each state, use the formal recipient
$V((z_1))\cdots((z_n))$ prescribed by
$|z_1|>\cdots>|z_n|$. Then
\begin{equation}\label{eq:vol2-wick-evaluator}
 Y(F_1,z_1)\cdots Y(F_n,z_n)
 ={:}\exp\left(\sum_{i<j}\Gamma_{ij}\right)
                       \prod_iF_i(j^{(i)}){:}.
\end{equation}
The exponential is finite on each list of polynomial inputs. Before
Laurent expansion, every coefficient is meromorphic with a finite pole
order on each pair diagonal. The formula commutes with differentiation
in the positions and with insertion of the vacuum. Grouping successive
inputs gives the same operator product in a common iterated Laurent
recipient.
\end{proposition}

\begin{proof}
Move each annihilation operator past the creation operators to its
right. Their commutators are scalars, so a term is specified by a set
of disjoint pairs of field slots belonging to different inputs. Each
pair contributes \eqref{eq:vol2-jet-contraction}; the unpaired slots
remain normally ordered. A product of $q$ contraction operators counts
the same unordered set of pairs $q!$ times, cancelled by the
exponential coefficient. Repeated derivatives count the choices among
identical slots. This proves the formula on monomials and hence on
polynomials.

Each contraction removes two slots. There are only finitely many
slots, so the exponential terminates, and each scalar coefficient is
a finite sum of products of the displayed rational functions.
Differentiating the contraction identity differentiates the corresponding
slot; the product rule accounts for all choices. A vacuum slot has no
jet variables, so it has no contractions and contributes the identity.
Finally, the operators $\Gamma_{ij}$ commute. Grouping the factors
partitions the same set of pairs into internal and cross pairs, each
used once. Their exponentials multiply to the single exponential in
\eqref{eq:vol2-wick-evaluator}. Expanding every rational coefficient in
one common order proves the last assertion.
\end{proof}

The meromorphic expression and its Laurent expansions have different
roles. At a collision with separation $u$, the differential-module
operation applied to a Laurent test $f(u)$ is
$\operatorname{pp}_u(f(u)Y(F,u)G)$.
A pole in the test can detect a regular field term. For example,
$\operatorname{pp}_u(u^{-1}Y(1,u)G)=u^{-1}G$.
Thus a vacuum insertion remains visible even though its field has no
singular part before multiplication by the test.

\begin{proposition}[All current modes]
\label{prop:vol2-polynomial-current-modes}
For every polynomial state $F$ and integer $m$,
\begin{equation}\label{eq:vol2-polynomial-current-modes}
 [J_m,Y(F,w)]
 =\lambda\sum_{r\geq0}(m)_{r+1}w^{m-r-1}
                 Y(\partial_{j_r}F,w),
 \qquad (m)_{r+1}=m(m-1)\cdots(m-r).
\end{equation}
The sum is finite, also for negative $m$.
\end{proposition}

\begin{proof}
A single current can contract with only one slot of $F$.
The contraction with $j_r(w)$ is
$\lambda(r+1)!(z-w)^{-r-2}$.
The coefficient of $(z-w)^{r+1}$ in the formal expansion of
$z^m$ about $w$ is $\binom m{r+1}w^{m-r-1}$, for every integer
$m$. Taking the residue in $z$ therefore gives the displayed falling
power. Only the finitely many jet variables occurring in $F$ contribute.
\end{proof}

Write $\langle-\rangle$ for the Gaussian vacuum functional, defined on a product of fields by
applying the product to $1\in V$ and then taking its constant
polynomial coefficient. In \eqref{eq:vol2-wick-evaluator} this retains
exactly the complete pairings. For fields $A_1,\ldots,A_n$, their
connected vacuum coefficient is
\[
 \kappa_n^{\mathrm{vac}}(A_1,\ldots,A_n)
 =\sum_{\pi\in\operatorname{Part}([n])}
   (-1)^{|\pi|-1}(|\pi|-1)!
   \prod_{B\in\pi}\left\langle\prod_{i\in B}A_i\right\rangle.
\]
Products within a block use the inherited order. This definition
specifies an algebraic vacuum functional; a numerical state or a
field-theoretic preparation requires its own construction.

Put $Q(z)={:}J(z)^2{:}$ and $K_{ij}=\lambda/(z_i-z_j)^2$.

\begin{proposition}[Connected quadratic fields]
\label{prop:vol2-wick-cycles}
For $n\geq3$,
\begin{equation}\label{eq:vol2-wick-cycles}
 \kappa_n^{\mathrm{vac}}(Q(z_1),\ldots,Q(z_n))
 =2^{n-1}\sum_{\sigma\in S_{\{2,\ldots,n\}}}
   K_{1,\sigma(2)}K_{\sigma(2),\sigma(3)}\cdots K_{\sigma(n),1}.
\end{equation}
For $n=2$ it is $2K_{12}^2$. In particular,
\begin{equation}\label{eq:vol2-wick-triangle}
 \kappa_3^{\mathrm{vac}}(Q(z_1),Q(z_2),Q(z_3))
 =\frac{8\lambda^3}{z_{12}^2z_{23}^2z_{13}^2},
 \qquad z_{ij}=z_i-z_j.
\end{equation}
These identities hold as meromorphic coefficients.
\end{proposition}

\begin{proof}
The partition formula cancels every disconnected pairing. Indeed,
for a fixed graph with $c$ connected components, its coefficient is
the sum of $(-1)^{|\pi|-1}(|\pi|-1)!$ over partitions of those
components. This sum is $1$ for $c=1$ and $0$ for $c>1$, as follows
by comparing coefficients in $\log(\exp t)=t$.

Normal ordering forbids a contraction within one input. A connected
pairing of two slots at each of at least three vertices is therefore
a cycle. Each unoriented cycle has $2^n$ labelled matchings, since
each vertex can assign its two slots to its neighbours in two ways.
Fixing vertex $1$ and permuting the other vertices counts that cycle
twice, once in each orientation. This gives the factor $2^{n-1}$.
At two vertices there are two bijections between the slots. At three
vertices there is one unoriented cycle and eight labelled matchings,
giving \eqref{eq:vol2-wick-triangle}.
\end{proof}

This triangle has no interaction vertex. It is already a correlation
of composite fields in a free theory. Its nonzero coefficient must be
retained by any comparison covering all polynomial boundary fields.
A calculation restricted to forests does not contain it.

\subsection{Changing the covariance}

The normal-order convention can itself be changed algebraically.
Let $B$ be a commutative $\mathbb Q$-algebra and
$A=B[q_1,\ldots,q_N]$. For symmetric matrices $K,R$ over $B$,
independent of the variables $q_i$, set
\[
 \Gamma_K=\sum_{i,j}K_{ij}\partial_{q_i}\otimes\partial_{q_j},
 \qquad m_K=m\exp(\Gamma_K),\qquad
 \Delta_R=\sum_{i,j}R_{ij}\partial_{q_i}\partial_{q_j},
 \qquad \alpha_R=\exp(\Delta_R/2),
\]
where $m$ is ordinary multiplication. All exponentials terminate on
polynomials. The variables may include separate copies at several
positions; the output operator then includes the cross derivatives
between those copies.

\begin{proposition}\label{prop:vol2-covariance-intertwiner}
For another symmetric constant matrix $S$,
\[
 \alpha_Rm_K=m_{K+R}(\alpha_R\otimes\alpha_R),
 \qquad \alpha_R\alpha_S=\alpha_{R+S}.
\]
\end{proposition}

\begin{proof}
The second-order product rule is
\[
 \tfrac12\Delta_Rm
 =m\bigl(\tfrac12\Delta_R\otimes1+
          1\otimes\tfrac12\Delta_R+\Gamma_R\bigr).
\]
All operators on the right commute. Exponentiating and composing
with $\exp(\Gamma_K)$ proves the first formula. The second follows
from $\Delta_R+\Delta_S=\Delta_{R+S}$.
\end{proof}

An intertwiner of these products does not specify a change of the
vacuum functional. Comparing states requires the functional as well
as the multiplication.

\section{Residue attachments and closed coincidences}
\label{sec:vol2-triangle-attachments}

The triangle has three dependent differences. Nevertheless its
meromorphic coefficient admits an explicit current extension using
only products in independent normal coordinates.
Work over $\mathbb C$, put $u=z_{12}$ and $v=z_{23}$, and use
the distributions $P_j$ of Proposition~\ref{prop:vol2-residue-current}.
Write $R_j=\bar\partial P_j$. Thus
\[
 P_j=\frac{(-1)^{j-1}}{(j-1)!}\partial_u^{j-1}P_1,
 \qquad
 R_j=\frac{\pi(-1)^{j-1}}{(j-1)!}
              \partial_u^{j-1}\delta_0\,d\bar u.
\]

The rational identity
\begin{equation}\label{eq:vol2-triangle-fractions}
 \frac1{u^2v^2(u+v)^2}
 =\frac1{v^4u^2}+\frac1{v^4(u+v)^2}
   -\frac2{v^5u}+\frac2{v^5(u+v)}
\end{equation}
follows by multiplying by the common denominator. It defines the
following extension of the triangle coefficient:
\begin{equation}\label{eq:vol2-triangle-current}
 \mathcal P_\triangle
 =P_4(v)\bigl(P_2(u)+P_2(u+v)\bigr)
       -2P_5(v)\bigl(P_1(u)-P_1(u+v)\bigr).
\end{equation}
Each summand is an external product in either $(u,v)$ or $(u+v,v)$.
Both pairs are linear coordinates, so every product in this formula
is defined.

Every meromorphic Wick coefficient also has a uniform principal-value
realization. For $F=G/Q$ on a smooth complex chart, it is
$[F]=\left.|Q|^{2s}G/Q\right|_{s=0}$, defined by analytic
continuation. This assignment is linear, independent of the
holomorphic presentation, and compatible with holomorphic derivatives;
see Andersson--Wulcan (2018)~\cite[\S2, Example~4.6, Corollary~4.26]{AnderssonWulcan2018}.
It applies to the finite rational coefficients of
Proposition~\ref{prop:vol2-wick-evaluator}, including cycles and
repeated edges. In independent coordinates it agrees with the
external products above, so $\mathcal P_\triangle$ is the
principal value of \eqref{eq:vol2-triangle-fractions}.
Restriction off the poles makes the assignment injective.
On its image, $[F]\diamond[G]=[FG]$ transports associative
meromorphic multiplication. This specified product does not extend
to arbitrary supported currents.

\begin{proposition}\label{prop:vol2-triangle-current}
The current $8\lambda^3\mathcal P_\triangle$ extends
\eqref{eq:vol2-wick-triangle}. Its complete Dolbeault attachment is
\begin{align}\label{eq:vol2-triangle-attachment}
 \bar\partial\mathcal P_\triangle
 &=R_4(v)\bigl(P_2(u)+P_2(u+v)\bigr)
   +P_4(v)\bigl(R_2(u)+R_2(u+v)\bigr)\notag\\
 &\quad-2R_5(v)\bigl(P_1(u)-P_1(u+v)\bigr)
   -2P_5(v)\bigl(R_1(u)-R_1(u+v)\bigr).
\end{align}
Holomorphic differentiation gives the attachments of its jets.
\end{proposition}

\begin{proof}
Away from the three divisors, the principal values are their
ordinary rational functions, so \eqref{eq:vol2-triangle-fractions}
proves the extension assertion. The distributional differential
satisfies the product rule on external products. Apply it separately
in the two coordinate systems to obtain all eight terms in
\eqref{eq:vol2-triangle-attachment}. Holomorphic derivatives commute
with $\bar\partial$ on distributions.
\end{proof}

The displayed extension fixes its contact terms. The rational
identity by itself would not exclude adding a current supported on
the collision. Nor does this construction define the ordinary product
of arbitrary singular distributions. The localization and contact
objects remain separate, as required by
Section~\ref{sec:vol2-supported-tensors}.

\subsection{The finite Koszul complex of a graph}

A closed coincidence gives a different complex. Let $\Gamma$ be
a finite oriented multigraph with vertex set $I$ and edge set $E$;
loops and repeated edges are allowed. Put
\begin{equation}\label{eq:vol2-graph-koszul}
 S=k[z_i:i\in I],\qquad
 K_\Gamma=S\otimes\Lambda(e_a:a\in E),\qquad
 |e_a|=-1,\quad de_{(i,j)}=z_i-z_j.
\end{equation}
No difference is inverted. A loop has zero differential.

Choose a spanning forest $T$ and one root in each component.
For each nonforest edge $a$, express its difference as a signed
sum along the unique forest path and set
\[
 c_a=e_a-\sum_{t\in\operatorname{path}(a)}\epsilon_t e_t.
\]
For a loop the sum is empty. The vectors $c_a$ form a basis of the
cycle space of the graph and satisfy $dc_a=0$.
Root coordinates and the forest differences $w_t$ form a linear
coordinate system: removing a leaf successively reconstructs every
vertex coordinate. With $B=k[z_{\mathrm{roots}}]$, the complex is
therefore
\[
 K_\Gamma=B[w_t:t\in T]\otimes\Lambda(e_t:t\in T)
                          \otimes\Lambda(c_a:a\notin T),
 \qquad d=\sum_tw_t\partial_{e_t}.
\]
Let $p$ set all $w_t,e_t$ to zero, and let $i$ include the resulting
root-and-cycle algebra.

\begin{proposition}[Forest contraction]
\label{prop:vol2-graph-contraction}
Let $N$ count total polynomial degree in the $w_t$ plus exterior
degree in the $e_t$, excluding the $c_a$. Define
\[
 H=\frac1N\sum_te_t\partial_{w_t}\quad(N>0),
 \qquad H=0\quad(N=0).
\]
Then
\[
 dH+Hd=1-ip,\qquad H^2=Hi=pH=0,
 \qquad H^*(K_\Gamma)=B\otimes\Lambda(c_a:a\notin T).
\]
The homotopy and the chosen section are $B$-linear. The augmentation
$p$ is $S$-linear for the quotient action on its target.
\end{proposition}

\begin{proof}
The graded commutator of $\sum_tw_t\partial_{e_t}$ and
$\sum_te_t\partial_{w_t}$ is the Euler operator in exactly the
$w_t,e_t$ variables. Both preserve $N$. Dividing by $N$ proves the
contraction identity on positive degree; on degree zero $ip$ is the
identity. The numerator of $H$ squares to zero because the exterior
multiplications anticommute and the ordinary derivatives commute.
Evaluation gives $pH=0$, and $Hi=0$ follows from the absence of
forest variables in the image of $i$. The target has zero
differential, proving the cohomology formula.

The formulas leave root coordinates unchanged. In contrast,
$H(w_t)=e_t$ whereas $w_tH(1)=0$, so $H$ is generally not
$S$-linear. The augmentation respects $S$ because setting all
forest differences to zero is precisely its quotient action.
\end{proof}

For the triangle the surviving generator is
$e_{13}-e_{12}-e_{23}$. Identifying the endpoints of the path
$12,23$ instead gives the closed generator $e_{12}+e_{23}$.
Keeping the original edge generators makes vertex identifications
maps of Koszul complexes. The choice of spanning forest calculates
their cohomology without removing those cycles.

The stable support complex of
Lemma~\ref{lem:vol2-redundant-support} is unchanged by a redundant
support equation. The finite Koszul complex retains its excess
cycle. Neither construction identifies that cycle with the
meromorphic triangle current. A graph loop, for example, gives a
closed exterior generator in \eqref{eq:vol2-graph-koszul}, but no
rational function $1/(z_i-z_i)$.

\section{Mixed kernels and their forest homotopies}
\label{sec:vol2-mixed-forest-homotopies}

To compare a mixed propagator with a boundary coefficient, its
point source must remain in the differential. Work on
$\mathbb R_t\times\mathbb C_z$, where $z=x+iy$, with orientation
$dt\wedge dx\wedge dy$. Partial forms are generated by $dt,d\bar z$,
and
\[
 D=dt\,\partial_t+d\bar z\,\partial_{\bar z},\qquad
 \partial_{\bar z}=\tfrac12(\partial_x+i\partial_y).
\]
Distribution coefficients are measured against $dt\,dx\,dy$.
Put $r=(t^2+|z|^2)^{1/2}$ and define
\begin{equation}\label{eq:vol2-mixed-axial-kernels}
 P=\frac{t\,d\bar z-2\bar z\,dt}{4\pi r^3},\qquad
 A=-\frac{\delta(t)\,dt}{\pi z},\qquad
 \Delta=\delta^{(3)}(t,z)\,dt\wedge d\bar z.
\end{equation}
The coefficients of $P$ are locally integrable. The coefficient of
$A$ is the external product of the one-dimensional delta and the
locally integrable Cauchy kernel.

\begin{lemma}\label{lem:vol2-kernel-source}
One has $DP=DA=\Delta$ as currents.
\end{lemma}

\begin{proof}
The coefficient of $dt\wedge d\bar z$ in $DP$ is
\[
 \frac1{4\pi}\left(\partial_t\frac{t}{r^3}
               +2\partial_{\bar z}\frac{\bar z}{r^3}\right).
\]
The imaginary angular terms cancel. The remainder is $1/(4\pi)$
times the divergence of $(t,x,y)/r^3$, whose flux through a sphere
is $4\pi$. It is therefore $\delta^{(3)}$ distributionally.
For $A$, use $\partial_{\bar z}(1/(\pi z))=\delta^{(2)}(z)$.
The minus sign in $A$ cancels the sign from moving $d\bar z$
past $dt$.
\end{proof}

Let $H(t)$ be the Heaviside distribution. Its value at zero is
irrelevant to the current. Define
\begin{equation}\label{eq:vol2-kernel-primitive}
 S(t,z)=\frac{H(t)-\tfrac12-\tfrac{t}{2r}}{\pi z}.
\end{equation}

\begin{proposition}\label{prop:vol2-kernel-primitive}
The function $S$ is locally integrable and satisfies
\[
 DS=P-A,\qquad
 \int_{\mathbb R}P=\int_{\mathbb R}A=-\frac1{\pi z}.
\]
At fixed $z$, its two tails are $O(|\bar z|/t^2)$.
\end{proposition}

\begin{proof}
On $t>0$ and $t<0$ respectively,
\[
 S_+=\frac{1-t/r}{2\pi z},\qquad
 S_-=\frac{-1-t/r}{2\pi z}.
\]
Both have derivatives
$\partial_tS_\pm=-\bar z/(2\pi r^3)$ and
$\partial_{\bar z}S_\pm=t/(4\pi r^3)$.
Their jump across $t=0$ is $1/(\pi z)$, supplying the
$\delta(t)/(\pi z)$ term in $P-A$. For fixed $t\ne0$ the
numerator vanishes at $z=0$, so there is no further axial delta.
The function is homogeneous of degree $-1$ with bounded angular
factor near the origin. Its first derivatives have degree $-2$.
A distribution supported only at the origin is a finite sum of
delta derivatives, whose degrees are at most $-3$. Thus no
additional point-supported term occurs, and the identity holds
on the whole space.

Integration of the $dt$ coefficient gives
\[
 -\frac{\bar z}{2\pi}\int_{-\infty}^{\infty}
          (t^2+|z|^2)^{-3/2}\,dt=-\frac1{\pi z}.
\]
Its absolute integral is locally integrable in $z$, which justifies
the distributional pushforward. The axial integral has the same
value. Expanding $t/r$ at either end gives the tail estimate.
\end{proof}

Multiplication by $-\pi$ converts this pushforward to the Cauchy
coefficient $1/z$, whose contour residue is one under
$(2\pi i)^{-1}\oint dz$. A physical coupling or pairing must
include that normalization. Every fixed finite holomorphic or
normal derivative commutes with the current identity.

\subsection{Cutoffs and boundary conditions}

For a smooth cutoff $\chi$ equal to one near the origin,
\begin{equation}\label{eq:vol2-kernel-cutoff}
 D(\chi P)=\Delta+D\chi\wedge P,\qquad
 D(\chi S)=\chi(P-A)+D\chi\,S.
\end{equation}
The annular terms belong to the supported comparison. If a
composition needs a proper projection, choose cutoffs giving that
property and retain these terms in its localization complex.

A smooth regulator follows from
$D^*=-\partial_t\iota_{dt}-4\partial_z\iota_{d\bar z}$.
For $0<\varepsilon<L$, set
\[
 K_\ell=(4\pi\ell)^{-3/2}e^{-r^2/(4\ell)},\qquad
 P_{\varepsilon,L}
 =D^*\left(\int_\varepsilon^L K_\ell\,d\ell
                           \;dt\wedge d\bar z\right).
\]
Since $DD^*+D^*D$ is the negative Euclidean Laplacian,
the heat equation gives
\[
 DP_{\varepsilon,L}=(K_\varepsilon-K_L)\,dt\wedge d\bar z.
\]
Both heat endpoints occur in its differential.

On a half-line the relative de Rham complex has Dirichlet
zero-forms and Neumann one-forms. Writing
$G_\ell(a)=(4\pi\ell)^{-1/2}e^{-a^2/(4\ell)}$, their kernels are
\[
 K^0_\ell(t,s)=G_\ell(t-s)-G_\ell(t+s),\qquad
 K^1_\ell(t,s)=G_\ell(t-s)+G_\ell(t+s).
\]
Differentiation gives $\partial_tK^0_\ell=-\partial_sK^1_\ell$.
Integration by parts, using the relative boundary value, proves
cochain compatibility. Tensoring with the tangential Dolbeault
complex preserves the identity. The polarization of a particular
field theory determines whether its normal complex is relative
or absolute.

\subsection{Forest cubes and contact terms}

Let $F$ be an oriented forest on a finite labelled set. For
$e=(i,j)$ put $t_e=t_i-t_j$ and $z_e=z_i-z_j$.
Choose one root per connected component. The root coordinates
and edge differences are a linear coordinate system with
determinant $\pm1$, by successive leaf deletion. Hence the
pullbacks $P_e,A_e,S_e,\Delta_e$ for different forest edges
have independent normal variables. Their ordered products,
including any fixed finite normal derivatives, are external
products of distributions in these coordinates.

\begin{theorem}\label{thm:vol2-forest-cubes}
On every labelled forest, the mixed and axial decorations are
joined by homotopy cubes compatible with contact contraction,
spectators, and products or graftings that remain forests.
This includes the edge $13$ with spectator $2$ and the crossed
pair of edges $13,24$.
\end{theorem}

\begin{proof}
Introduce $s_e\in[0,1]$ for each edge, with parameter forms
written before spatial forms, and set
\[
 \mathbb P_e=A_e+s_e(P_e-A_e)+ds_e\,S_e.
\]
Then $(D+d_s)\mathbb P_e=\Delta_e$: the term
$ds_e(P_e-A_e)$ cancels $-ds_eDS_e$.
Independence of the edge variables defines every forest product.
Its differential is the signed sum replacing one edge by its
contact current. At $s_e=0$ or $1$ the edge becomes $A_e$ or
$P_e$. Restrictions to two faces commute, giving the higher
face identities. The diagonal interval varies all edges together.

Contact contraction identifies the two endpoints of an edge;
the remaining differences are coordinates of the contracted
forest. The distributional equations persist under finite normal
differentiation and retain its contact jets. Adding an isolated
vertex adds a spectator, and the empty product is the coefficient
unit. Edge reorderings have the usual form signs.
\end{proof}

For two edges define $H_{e,f}=S_eP_f-A_eS_f$. Since
$|P|=|A|=1$ and $|S|=0$, the product rule gives
\begin{equation}\label{eq:vol2-two-edge-contact}
 P_eP_f-A_eA_f
 =DH_{e,f}+\Delta_eS_f-S_e\Delta_f.
\end{equation}
The last two terms express the common point sources of the
kernels. They remain even though $P-A$ is exact.

At three positions put $z_{12}=u$, $z_{23}=v$, $z_{13}=u+v$.
The axial products have the same time delta and oriented time
two-form. Their scalar coefficients satisfy
\[
 \frac1{uv}-\frac1{u(u+v)}-\frac1{v(u+v)}=0.
\]
Each term is locally integrable on $\mathbb C^2$ in its
independent linear coordinates. Thus the rational identity is
also a distributional identity. Define
$\mathcal K=H_{12,23}-H_{12,13}-H_{13,23}$. Applying
\eqref{eq:vol2-two-edge-contact} to the three pairs gives
\begin{align}\label{eq:vol2-three-channel-current}
 P_{12}P_{23}-P_{12}P_{13}-P_{13}P_{23}
 &=D\mathcal K+\mathcal C,\notag\\
 \mathcal C
 &=\Delta_{12}S_{23}-S_{12}\Delta_{23}
   -\Delta_{12}S_{13}+S_{12}\Delta_{13}\notag\\
 &\quad-\Delta_{13}S_{23}+S_{13}\Delta_{23}.
\end{align}
Every product here has independent normal variables.

These are scalar coefficient currents on the mixed space.
Decorating them by fields requires a map preserving their
contact terms, translations, partial vacua and supported
attachments. A cycle has dependent edge differences, so its
mixed current product needs an additional extension. The
meromorphic triangle of \eqref{eq:vol2-triangle-current} supplies
a specified boundary coefficient, without identifying it with
such a mixed product.

\section{Classes discarded by consecutive cuts}
\label{sec:vol2-discarded-collisions}

Consecutive cuts preserve an order on the labels. A collision can
join labels that are separated in that order. To measure this
difference, we compare two coefficient complexes before introducing
a map from coefficients to fields.

Let $R_n=k[z_1,\ldots,z_n]$ and
$L_n=R_n[(\prod_{i<j}z_{ij})^{-1}]$.
An ordered partition $P=(B_1|\cdots|B_r)$ has coefficient module
\begin{equation}\label{eq:vol2-partition-coefficient}
 M_P=R_n[z_{ij}^{-1}:i,j\text{ lie in different blocks of }P]
       \subset L_n.
\end{equation}
Equal subrings attached to different ordered partitions remain
different summands. Tensor every term with the common right
density $\omega_n=\omega_{\mathbb A^n}$; for a constant density
$\nu$, $(f\nu)\partial_i=-(\partial_i f)\nu$.
We suppress this common factor in the coordinate formulas.

Write $e_P=\xi_{B_1}\cdots\xi_{B_r}$, where
$|\xi_B|=1-|B|$. The full complex $C_n^{\mathrm{all}}$ has
term $M_Pe_P$ in degree $r-n$. Its differential splits a block
$B_i=A\sqcup C$ into two nonempty ordered parts, with coefficient
the localization inclusion and sign
\begin{equation}\label{eq:vol2-partition-sign}
 (-1)^{\sum_{j<i}(1-|B_j|)+|A|+1}\epsilon(A,C).
\end{equation}
Here $\epsilon(A,C)$ is the sign of the concatenated increasing
lists of $A,C$ relative to the increasing list of $B_i$.
A partition is \emph{consecutive} when its blocks are intervals
in increasing order. Projection to those summands gives the
consecutive-cut complex $C_n^{\mathrm{cut}}$. Its kernel is
\begin{equation}\label{eq:vol2-discarded-complex}
 F_n^{r-n}=\bigoplus_{\substack{P\text{ has }r\text{ blocks}\\
                                  P\text{ is not consecutive}}}M_Pe_P.
\end{equation}
There is no one-block term in $F_n$.

\begin{lemma}\label{lem:vol2-partition-complex}
The discarded summands form a subcomplex and $d^2=0$.
After extending every coefficient to $L_n$, the full complex is
the oriented cellular chain complex of the permutahedron of
dimension $n-1$, with chain degree negated. For $n\geq3$, the
map from discarded two-block terms to three-block terms is
injective, both after localization and on the actual modules.
\end{lemma}

\begin{proof}
If a refinement is consecutive, merging its two adjacent new
blocks recovers a consecutive partition. Thus a nonconsecutive
partition has only nonconsecutive refinements. Two splits in
different blocks cancel by the odd Leibniz sign. Two successive
splits of the same block have the same density permutation and
opposite cobar signs. Hence $d^2=0$.

The permutahedron is the convex hull of the permutations of
$(1,\ldots,n)$. A face is specified by an ordered partition;
its dimension is $n-r$, and splitting one block gives a boundary
face. Orient the block factor by the increasing labels, then
orient the product in block order. Its boundary orientation is
exactly \eqref{eq:vol2-partition-sign}. This identifies the
fully localized incidence complex.

The permutahedron is a ball. In its cellular complex the kernel
on facets is the image of the unique top cell, spanned by
$d\xi_{[n]}$. Every facet has coefficient $1$ or $-1$ in that
boundary. The discarded facet space omits all $n-1$ consecutive
facets, so it contains no nonzero multiple of the top boundary.
This proves injectivity over $k$, hence over $L_n$. Each actual
coefficient module embeds in $L_n$, giving the last assertion.
\end{proof}

\subsection{The outer principal-part module}

Put $u=z_{1n}$. Retain in $d\xi_{[n]}$ only those facets which
keep $1,n$ together:
\begin{equation}\label{eq:vol2-outer-numerator}
 T_n=\sum_{\substack{A\sqcup C=[n],\ A,C\ne\varnothing\\
                        1,n\text{ lie in the same part}}}
          (-1)^{|A|+1}\epsilon(A,C)\xi_A\xi_C,
 \qquad b_n=dT_n.
\end{equation}
These facets are all discarded. The degrees of $T_n$ and $b_n$
are $2-n$ and $3-n$, respectively.

\begin{theorem}[Outer classes in every arity]
\label{thm:vol2-outer-injection}
For every $n\geq3$, the formula
\begin{equation}\label{eq:vol2-outer-injection}
 \omega_n\otimes_{R_n}(R_n[u^{-1}]/R_n)
       \lhook\joinrel\longrightarrow H^{3-n}(F_n),
 \qquad [f]\longmapsto[fb_n]
\end{equation}
defines an injective map of right differential modules.
Thus no arity-independent lower bound exists for the degrees
of nonzero cohomology of the discarded complexes.
\end{theorem}

\begin{proof}
Consider a three-block partition which still keeps $1,n$
together. Its two facet coarsenings both occur in $T_n$, so
their contributions to $dT_n$ cancel. Every nonzero coordinate
of $b_n$ therefore separates $1,n$ and has a coefficient module
in which $u$ is invertible. Consequently $fb_n$ is an allowed
cochain for every $f\in R_n[u^{-1}]$, and it is closed.
If $f\in R_n$, the primitive $fT_n$ is allowed, so the formula
factors through principal parts.

Conversely, suppose $fb_n=dv$ for an actual discarded facet
cochain $v$. Over $L_n$, the injectivity in
Lemma~\ref{lem:vol2-partition-complex} forces $v=fT_n$.
Choose a facet $P$ occurring in $T_n$. Its module does not
invert $u$, and unique factorization gives
$M_P\cap R_n[u^{-1}]=R_n$. The corresponding coordinate of
$v$ is $\pm f$, so $f$ is regular. Such a facet exists for
$n\geq3$, and the same injectivity shows $b_n\ne0$.

The coefficients of $b_n$ are constants. Differentiating
$fb_n$ therefore differentiates only $f$. Restoring the common
right density proves differential-module linearity. Every normal
derivative of the simple principal part is retained.
\end{proof}

For four positions the numerator is
\begin{equation}\label{eq:vol2-outer-four-numerator}
 \begin{split}
 T_4={}&-\xi_{124}\xi_3+\xi_{134}\xi_2-\xi_{14}\xi_{23}\\
       &-\xi_2\xi_{134}-\xi_{23}\xi_{14}+\xi_3\xi_{124}.
 \end{split}
\end{equation}
Thus $z_{14}^{-1}dT_4$ is a nonzero class in $H^{-1}(F_4)$.
Its fully localized primitive $z_{14}^{-1}T_4$ has a forbidden
pole in each displayed facet coefficient. The obstruction is
precisely that failure to lie in the specified module.

The relation with the cut class also fixes a sign and a degree.
The degreewise exact sequence is
\[
 0\longrightarrow F_n\longrightarrow C_n^{\mathrm{all}}
                  \longrightarrow C_n^{\mathrm{cut}}\longrightarrow0.
\]
Its cut top boundary is
\[
 U_n=d_{\mathrm{cut}}\xi_{[n]}
    =\sum_{p=1}^{n-1}(-1)^{p+1}\xi_{[1,p]}\xi_{[p+1,n]}.
\]
Every term separates $1,n$, so $fU_n$ is an allowed cut
cocycle for every $f\in R_n[u^{-1}]$.

\begin{proposition}\label{prop:vol2-outer-connecting}
The connecting homomorphism sends
\[
 H^{2-n}(C_n^{\mathrm{cut}})\longrightarrow H^{3-n}(F_n),
 \qquad [fU_n]\longmapsto-[fb_n].
\]
\end{proposition}

\begin{proof}
Write $d\xi_{[n]}=U_n+T_n+V_n$, where $V_n$ consists of
discarded facets separating $1,n$. Lift $fU_n$ to the same
cochain in the full complex. Its differential is
$d(fU_n)=-fb_n-d(fV_n)$. The latter term is a legal boundary
in $F_n$, since every facet in $V_n$ inverts $u$.
\end{proof}

This comparison concerns restriction of all-subset cobar
coefficients to a chosen order and cobar after restriction to
cuts. Its kernel does not refute bar--cobar duality on either
of its stated categories. A map from these coefficient classes
to a field-action complex remains additional data.

\subsection{Global cohomology at four positions}

The full four-position permutahedron has $14$ facets, $36$
edges and $24$ vertices. Removing consecutive partitions removes
three facets, three edges and the ascending vertex. Hence
\begin{equation}\label{eq:vol2-four-discarded-matrices}
 F_4:\qquad \bigoplus_{11}M_P\xrightarrow{A}
                 \bigoplus_{33}M_Q\xrightarrow{B}L^{23},
 \qquad L=L_4,
\end{equation}
in degrees $-2,-1,0$. The fully localized ranks are $11,22$.
For the second rank, the edges connect all remaining vertices:
adjacent swaps increasing inversion number connect every
nonascending permutation to the reverse permutation without
passing through the removed vertex. Thus over $L$,
$\ker B=\operatorname{im}A$ and $\ker A=0$.

For an unordered pair $h$, write
\[
 L_h=R_4[z_{ab}^{-1}:\{a,b\}\ne h],\qquad W_h=L/L_h.
\]
A three-block term has a unique unsplit pair $h$ and coefficient
$L_h$. Let $I_h$ be the discarded facets keeping $h$ together
and $J_h$ the three-block terms keeping it together. Denote the
corresponding submatrix of $A$ by $K_h$.

\begin{lemma}\label{lem:vol2-four-pair-ranks}
The sizes and ranks are
\[
 \begin{array}{c|rrrrrr}
 h&12&13&14&23&24&34\\\hline
 |I_h|&4&5&6&4&5&4\\
 |J_h|&5&6&6&5&6&5\\
 \operatorname{rank}K_h&4&5&5&4&5&4.
 \end{array}
\]
Only $K_{14}$ has a kernel. In the order of the six terms in
\eqref{eq:vol2-outer-four-numerator}, that kernel is generated by
$(-1,1,-1,-1,-1,1)$.
\end{lemma}

\begin{proof}
Use \eqref{eq:vol2-partition-sign}. In the displayed order,
\[
 K_{14}=\begin{pmatrix}
 1&0&-1&0&0&0\\0&1&1&0&0&0\\-1&0&0&1&0&0\\
 0&0&0&-1&1&0\\0&-1&0&0&0&1\\0&0&0&0&-1&-1
 \end{pmatrix}.
\]
Its equations determine every coordinate from the first. For
the other pairs, order the partitions lexicographically by
their increasing block lists. The rows for $12,23,34$ give
respectively
\[
 \begin{array}{c|ccccc}
 12&-v_1&-v_2&v_2+v_3&v_1-v_4&-v_3+v_4\\
 23&v_1&-v_2&v_2-v_3&-v_1+v_4&v_3-v_4\\
 34&v_1&v_2&-v_2+v_3&-v_1-v_4&-v_3+v_4.
 \end{array}
\]
Each zero system forces all four coordinates to vanish. For
$13$ the rows are
$-v_1,v_1-v_2,-v_3,v_3+v_4,v_2+v_5,-v_4-v_5$;
for $24$ they are
$v_1,v_2,-v_1-v_4,v_3+v_4,-v_2+v_5,-v_3-v_5$.
These force all five coordinates to vanish. Every elimination
uses only $1,-1$, so the same conclusions hold with coefficients
in $W_h$.
\end{proof}

\begin{theorem}\label{thm:vol2-four-outer-cohomology}
Globally on the standard affine chart,
\[
 H^{-2}(F_4)=0,\qquad
 H^{-1}(F_4)=\omega_4\otimes_{R_4}
                    (R_4[z_{14}^{-1}]/R_4).
\]
The second identification sends $[f]$ to $[fdT_4]$.
\end{theorem}

\begin{proof}
The first assertion follows from facet injectivity. Suppress
the density for the second. A degree-minus-one cocycle $w$ has
a unique fully localized primitive $v\in L^{11}$, with $Av=w$.
Requiring the $J_h$ coordinates of $w$ to lie in $L_h$ gives
\[
 K_h(v_i\bmod L_h)_{i\in I_h}=0.
\]
For $h\ne14$, Lemma~\ref{lem:vol2-four-pair-ranks} forces each
indicated coordinate of $v$ to be regular along $h$. Unique
factorization gives, inside $L$,
\[
 M_P=\bigcap_{h\text{ internal to }P}L_h.
\]
Putting $u=z_{14}$, we obtain $v_P\in M_P[u^{-1}]$.
If $P$ separates $1,4$, this is already $M_P$, and its
coordinate is a legal boundary. The six remaining principal
parts are proportional to the coefficients of $T_4$. Thus
$H^{-1}$ is the intersection, inside $W_{14}$, of the images
of $M_P[u^{-1}]/M_P$ for these six facets.

To compute that intersection at the other collision divisors
as well, use
\[
 u=z_{14},\quad a=z_{12},\quad b=z_{23},\quad
 c=a+b=z_{13},\quad y=z_4,
 \qquad D=k[y,a,b].
\]
The remaining differences are $z_{24}=u-a$, $z_{34}=u-c$.
Finite negative Laurent expansion at $u=0$ identifies
\[
 W_{14}=\bigoplus_{j\geq1}D[(abc)^{-1}]u^{-j}.
\]
Indeed, expand $(u-a)^{-1},(u-c)^{-1}$ only to the finite
order required by the pole in $u$. A zero negative part means
that the numerator is divisible by that power of $u$, hence
the original fraction belongs to $L_{14}$.

The three distinct facet images, from $124|3$, $134|2$ and
$14|23$, are respectively
\[
 \bigoplus_{j\geq1}D[(bc)^{-1}]u^{-j},\quad
 \bigoplus_{j\geq1}D[(ab)^{-1}]u^{-j},\quad
 \bigoplus_{j\geq1}D[(ac)^{-1}]u^{-j}.
\]
The reversed facets repeat them. For example, $M_{124|3}$
inverts $b,c,u-c$, whose negative Laurent coefficients invert
only $b,c$; all such coefficients occur. The irreducibles
$a,b,c=a+b$ are pairwise nonassociated in $D$, so the
intersection of the three coefficient rings is $D$.
The principal-part intersection is therefore
$\bigoplus_{j\geq1}Du^{-j}=R_4[u^{-1}]/R_4$.
The identification is $[f]\mapsto[fdT_4]$, already proved
differential-linear in Theorem~\ref{thm:vol2-outer-injection}.
\end{proof}

Degree zero has a finite global presentation. Let $G$ be the
permutahedron graph on $S_4\setminus\{1234\}$ and choose
$\sigma_0=4321$. Label each edge by its swapped pair $h$ and
define
\[
 \mathcal G_{R_4}=
 \frac{\displaystyle\bigoplus_{\sigma\ne\sigma_0}L e_\sigma}
      {\displaystyle\sum_{\{\sigma,\tau\}\in E(G)}
            L_{h(\sigma,\tau)}(e_\sigma-e_\tau)},
 \qquad e_{\sigma_0}=0.
\]
Then
\begin{equation}\label{eq:vol2-four-H0}
 H^0(F_4)=\omega_4\otimes_{R_4}(L\oplus\mathcal G_{R_4}).
\end{equation}
To see this, orient the singleton density basis so that $B$
is the graph incidence map. Its image lies in the kernel of
the sum of vertex coordinates. Inclusion at $\sigma_0$ splits
that sum. In its kernel, the basis
$e_\sigma-e_{\sigma_0}$ gives exactly the displayed relations.
All basis changes are constant and preserve the differential
module structure. This presents the global module by $22$
generator modules and $33$ relation modules; it does not
decompose its extensions across all diagonal intersections.

\subsection{The crossed four-position chart}

On the open set separating $\{1,3\}$ from $\{2,4\}$, put
\[
 \begin{gathered}
 u=z_{13},\quad v=z_{24},\quad
 A=R_4[(z_{12}z_{14}z_{23}z_{34})^{-1}],\\
 A_u=A[u^{-1}],\quad A_v=A[v^{-1}],\quad L=A[(uv)^{-1}].
 \end{gathered}
\]
The two remaining collision coordinates $u,v$ are independent.
The coefficient types in \eqref{eq:vol2-four-discarded-matrices}
become
\[
 A^2\oplus A_u^3\oplus A_v^3\oplus L^3
 \longrightarrow A_u^6\oplus A_v^6\oplus L^{21}
 \longrightarrow L^{23}.
\]

A nonzero scalar entry $a$ between two identical coefficient
modules can be cancelled. Isolating that entry by triangular
row and column operations gives a contractible two-term
summand with homotopy $a^{-1}$. The remaining differential is
the Schur complement
$D'_{rc}=D_{rc}-D_{rj}a^{-1}D_{ic}$.
Every entry still composes allowed localization inclusions.
These changes commute with differential operators.

Order the partitions lexicographically by increasing block
lists. Cancel the first available scalar between identical
modules, using lower degree first, then row, then column.
The nine cancellations from degree $-2$ have types
\[
 A_u,L,A_v,A_v,L,A_u,A_u,L,A_v
\]
and pivots $1,1,-1,-1,1,1,1,1,1$.
The next eighteen cancellations, from degree $-1$, have type
$L$ and pivot $1$. The remaining complex is
\[
 A^2\xrightarrow{M_0}A_u^3\oplus A_v^3
                              \xrightarrow{M_1}L^5.
\]
In its retained middle order $(A_u,A_v,A_v,A_u,A_u,A_v)$,
\[
 M_0=\begin{pmatrix}1&0\\-1&0\\0&1\\0&-1\\-1&1\\1&-1\end{pmatrix},
 \qquad
 M_1=\begin{pmatrix}
 1&1&0&0&0&0\\0&0&1&1&0&0\\0&-1&-1&0&1&0\\
 -1&0&0&-1&0&1\\0&0&0&0&-1&-1
 \end{pmatrix}.
\]
Use the basis matrices
\[
 S=\begin{pmatrix}
 1&0&1&0&0&0\\0&0&0&1&0&1\\0&0&0&0&-1&0\\
 0&-1&0&0&0&0\\-1&1&0&0&0&0\\0&0&0&-1&1&0
 \end{pmatrix},\qquad
 W=\begin{pmatrix}
 1&0&1&1&1\\0&-1&0&0&0\\-1&1&0&-1&0\\
 -1&1&-1&0&0\\1&-1&0&0&0
 \end{pmatrix}.
\]
Both are invertible. The first three columns of $S$ have
coefficient $A_u$, the last three $A_v$, so this basis change
respects the modules. Matrix multiplication gives
\begin{equation}\label{eq:vol2-crossed-canonical}
 S^{-1}M_0=
 \begin{pmatrix}1&0\\0&1\\0&0\\-1&0\\0&-1\\0&0\end{pmatrix},
 \qquad
 W^{-1}M_1S=
 \begin{pmatrix}
 1&0&0&1&0&0\\0&1&0&0&1&0\\0&0&1&0&0&0\\
 0&0&0&0&0&1\\0&0&0&0&0&0
 \end{pmatrix}.
\end{equation}

\begin{theorem}[Crossed local form]
\label{thm:vol2-crossed-collision}
On this open set, $F_4$ is chain-homotopy equivalent by right
differential-module maps to
\begin{equation}\label{eq:vol2-crossed-splitting}
 L[0]\oplus[A_u\to L]\oplus[A_v\to L]
       \oplus[A\to A_u\oplus A_v\to L]^{\oplus2}.
\end{equation}
The two-term complexes have degrees $-1,0$; the three-term
complexes have degrees $-2,-1,0$. In particular,
\[
 H^{<0}(F_4)=0,\qquad
 H^0(F_4)=\omega_4\otimes_A
 \left(L\oplus L/A_u\oplus L/A_v
              \oplus\bigl(L/(A_u+A_v)\bigr)^{\oplus2}\right).
\]
\end{theorem}

\begin{proof}
The first two columns and their fourth and fifth partners
in \eqref{eq:vol2-crossed-canonical} give two complexes with
maps $a\mapsto(a,-a)$ and $(f,g)\mapsto f+g$.
Columns three and six give the two remaining two-term
complexes. The last target row gives $L[0]$.
All preceding cancellations were chain homotopy equivalences.

The maps $A_u\to L$ and $A_v\to L$ are injective. In a
three-term summand the middle kernel is $(f,-f)$ with
$f\in A_u\cap A_v=A$. The last equality follows by unique
factorization in the localization of $R_4$: the independent
prime factors $u,v$ remain nonassociated and neither has been
inverted in $A$. Only the displayed cokernels remain.
\end{proof}

The codimension-two module $L/(A_u+A_v)$ is nonzero. Near
$u=v=0$, its class $1/(uv)$ has a negative power in both
independent normals, which cannot occur in a sum of terms
with poles in only one of them. The crossed chart has retained
these two copies of the collision module. It has inverted
$z_{14}$, so its vanishing negative cohomology agrees with the
global outer class, whose support is $z_{14}=0$.

At a generic point of that outer divisor, where all other
differences are invertible, set
$u=z_{14}$ and $A=R_4[(z_{12}z_{13}z_{23}z_{24}z_{34})^{-1}]$.
The same cancellations give
\[
 [A\to A[u^{-1}]]_{[-2,-1]}
 \oplus[A\xrightarrow{(1,-1)}A[u^{-1}]^2]_{[-1,0]}.
\]
Thus $H^{-1}=A[u^{-1}]/A$ and
$H^0=A[u^{-1}]\oplus A[u^{-1}]/A$ there. These are full
principal-part modules, including every normal derivative.

\section{Independent normals and the full diagonal}
\label{sec:vol2-full-diagonal-support}

The crossed calculation retains proper collision strata. The
full diagonal has a simpler normal calculation, whose shift
depends on the number of independent equations. We compute it
before imposing any identification with a field coefficient.

Let $X=\mathbb A^b_k$ have coordinates $y_1,\ldots,y_b$, let
$Y=X\times\mathbb A^c_k$ have normal coordinates
$u_1,\ldots,u_c$, and let $i:X\hookrightarrow Y$ be the zero
section. For a left $D_X$-module $L$, its normal direct image is
\[
 i_+L=L[\eta_1,\ldots,\eta_c],\qquad
 \partial_{u_j}=\eta_j,\qquad u_j=-\partial_{\eta_j}.
\]
Tangential operators act on $L$. All polynomial normal jets are
retained. Write $\operatorname{DR}_X(L)$ for the unshifted
algebraic de Rham complex in degrees $0,\ldots,b$.

\begin{proposition}\label{prop:vol2-independent-normal-support}
There is an explicit contraction
\[
 R\operatorname{Hom}_{D_Y}(\mathcal O_Y,i_+L)
       \simeq\operatorname{DR}_X(L)[-c].
\]
For right differential modules, put
$L^r=\omega_X\otimes_{\mathcal O_X}L$ and
$\mathbb J_Y=\omega_Y[b+c]$. Then
\begin{equation}\label{eq:vol2-independent-density-shift}
 R\operatorname{Hom}_{D_Y}
       \bigl(\mathbb J_Y,i_*(L^r[b])\bigr)
       \simeq\operatorname{DR}_X(L)[-2c],
\end{equation}
where $i_*$ denotes right differential-module direct image.
\end{proposition}

\begin{proof}
Apply Hom from the Spencer resolution of $\mathcal O_Y$.
One normal direction contributes the cochain complex
\[
 K_j=[k[\eta_j]\xrightarrow{\eta_j}k[\eta_j]]
\]
in degrees zero and one. Let $p_j$ extract the constant
coefficient in degree one and let $i_j$ insert it. Let $h_j$
vanish on degree zero and send $\eta_j^a$ in degree one to
$\eta_j^{a-1}$ for $a>0$, and $1$ to zero. Then
\[
 dh_j+h_jd=1-i_jp_j,\qquad p_ji_j=1,
 \qquad h_j^2=h_ji_j=p_jh_j=0.
\]
On the tensor product of the normal complexes set
\[
 H=\sum_{j=1}^c(i_1p_1)\otimes\cdots\otimes(i_{j-1}p_{j-1})
                         \otimes h_j\otimes1\otimes\cdots\otimes1,
\]
using the Koszul signs for tensor products of graded maps.
The identity follows by telescoping
$1-q_1\cdots q_c=\sum_jq_1\cdots q_{j-1}(1-q_j)$, with
$q_j=i_jp_j$. The cross terms in $H^2$ contain $h_ji_j$ or
$p_jh_j$ and vanish. Tensoring with the tangential de Rham
complex contracts the whole Spencer complex onto that complex
in normal degree $c$. Every polynomial jet has finite degree;
the argument also applies to an infinite-rank module $L$.

Left-to-right conversion identifies the source and target of
\eqref{eq:vol2-independent-density-shift} with
$\mathcal O_Y[b+c]$ and $i_+L[b]$. Their relative shift is
$-c$. Combining it with the normal Spencer degree gives
$-2c$.
\end{proof}

For an ordered $n$-point affine chart, the first coordinate
and its differences with the other points give one tangential
coordinate and $n-1$ independent normals. Let
$N=\operatorname{End}_{k[y]}(\mathcal V)$, with its commutator
connection, be the polynomial-root endomorphism module of
Theorem~\ref{thm:vol2-polynomial-operators}. Its de Rham
complex is $G(N)$, with cohomology $B_0$ in degree zero and
$W$ in degree one.

Suppose a specified collision comparison identifies the actual
full-diagonal coefficient of an ambient object with $N^r[1]$.
The supported linear mapping complex from
$\mathbb J_{\mathbb A^n}$ is then
\begin{equation}\label{eq:vol2-full-diagonal-groups}
 G(N)[-2(n-1)],\qquad
 H^{2n-2}=B_0,\quad H^{2n-1}=W,
\end{equation}
and all other cohomology vanishes. This follows from
Proposition~\ref{prop:vol2-independent-normal-support} with
$b=1,c=n-1$ and the stated endomorphism contraction. Replacing
the source by $\mathbb J_{\mathbb A^n}[-d]$ changes the complex
to $G(N)[d-2(n-1)]$.

The collision identification is a hypothesis about an actual
coefficient object and its maps. Matching a rank does not supply
it. For $n\geq3$, the complement of the full diagonal still
contains proper collision strata; their simultaneous support
requires their attachment maps. The localization triangle also
controls the map from supported to ambient cohomology, which
need not be injective. Finally, dependent graph equations have
codimension equal to their rank, and this calculation concerns
linear Hom. Coalgebra and substitution compatibility require
their respective operation maps.

\section{Heisenberg compositions on the pair support}
\label{sec:vol2-supported-heisenberg}

We now construct an outer operation on a complete pair-support
complex. Its source includes poles in the remaining separation
and the triple-diagonal term. A correction to the binary path
will make the operation closed on this larger domain.

Fix $\kappa\in k^\times$ and $X=\mathbb A^1_k$. The dg vertex
algebra is $\mathcal Z=V_\kappa\otimes\Lambda[c]$, with
$|J|=0$, $|c|=1$, $dJ=\kappa Tc$ and $dc=0$.
Its state complex is
\begin{equation}\label{eq:vol2-supported-state-complex}
 C=k[u_0,u_1,\ldots]\otimes\Lambda(c_0,c_1,\ldots),
 \quad u_n=T^nJ,\quad c_n=T^nc,\quad
 du_n=\kappa c_{n+1},\quad dc_n=0.
\end{equation}
All states and state tensor products are algebraic. The current
bracket is $[J_\nu J]=\kappa\nu\mathbf1$; brackets with $c$
vanish. The Fourier variable $\nu$ records normal derivatives:
the contact of $v^{-n-1}$ is the $n$th normal derivative divided
by $n!$. On regular state inputs this is the usual bracket
$[a_\nu b]=\operatorname{Res}_v e^{\nu v}Y(a,v)b\,dv$.

\subsection{A common resolution and its state contraction}

Let $E$ be the free quartet vertex algebra on fields of degrees
$(0,0,1,-1)$. Form $\mathcal W=V_\kappa\otimes E^{\otimes3}$,
with quartet generators $(t_s,p_s,\eta_s,\iota_s)$ and
\[
 Qt_s=\eta_s,\qquad Q\iota_s=p_s,
 \qquad [p_s{}_\nu t_l]=[\iota_s{}_\nu\eta_l]=\delta_{sl}.
\]
The other generating brackets vanish, apart from those forced
by graded skew symmetry and translation. In this ambient
algebra take $QJ=0$; its current commutes with the quartets.
Put
\[
 K_3=J+p_1+p_2+p_3,\qquad
 \mathcal B=\operatorname{Com}(K_3,\mathcal W),\qquad
 U_s=J+\sum_{l<s}p_l-\kappa Tt_s.
\]
The commutant means the common kernel of the nonnegative modes
of $K_3$. There are vertex chain maps
\begin{equation}\label{eq:vol2-supported-slots}
 j_s:\mathcal Z\longrightarrow\mathcal B,
 \qquad j_s(J)=U_s,\quad j_s(c)=-\eta_s.
\end{equation}
Distinct images $j_s(\mathcal Z)$ commute.
Indeed $QU_s=-\kappa T\eta_s$ and
$[U_r{}_\nu U_s]=\delta_{rs}\kappa\nu$.
For $r<s$, the mixed contribution
$[-\kappa Tt_r{}_\nu p_r]=-\kappa\nu$ cancels the
current bracket. All current--ghost and ghost--ghost
brackets are zero, and these images commute with $K_3$.
The assignments therefore extend to the stated vertex maps.

Here is a direct description of the state space of $\mathcal B$.
Write $\tau^{(s)}_n=T^nt_s/n!$ for the divided quartet jets,
with $n\geq0$, and use the current modes $x_n=J_{-n}1$.
In these polynomial coordinates, the nonnegative modes of $K_3$
are
\[
 \sum_s\partial_{\tau^{(s)}_0},\qquad
 \kappa n\partial_{x_n}+\sum_s\partial_{\tau^{(s)}_n}\quad(n>0).
\]
Changing coordinates to
$x_n-\kappa n\tau^{(1)}_n$ and
$\tau^{(s)}_n-\tau^{(1)}_n$ for $s=2,3$ leaves these as
derivatives in the remaining $\tau^{(1)}_n$ coordinates.
Their common kernel is therefore freely spanned by the PBW
monomials in the translates of
\[
 U_1,\eta_1;\qquad
 a_2,a_3,e_2,e_3;\qquad
 \iota_1,\iota_2,\iota_3,p_1,p_2,p_3,
 \quad a_s=t_s-t_1,\quad e_s=\eta_s-\eta_1.
\]
Their differentials are
$QU_1=-\kappa T\eta_1$, $Qa_s=e_s$, $Q\iota_s=p_s$,
with the remaining images zero. The fields after the first
semicolon are auxiliary.

Let $r:\mathcal B\to C$ set every auxiliary coordinate to
zero, send $U_1$ to $J$ and $\eta_1$ to $-c$. It is a map of
state complexes. Let $N$ count total auxiliary number and let
$s$ be the odd polynomial derivation with
\[
 s(T^me_j)=T^ma_j\quad(j=2,3),\qquad
 s(T^mp_j)=T^m\iota_j\quad(j=1,2,3),
\]
zero on the other generators. For $q$ tensor factors use the
graded sum $s_q$ and total number $N_q$, and put
$h_q^{\mathrm{aux}}=s_q/N_q$ on positive auxiliary number,
zero on number zero.

\begin{lemma}\label{lem:vol2-supported-state-contraction}
With $i=j_1$, these maps satisfy
\[
 dh_q^{\mathrm{aux}}+h_q^{\mathrm{aux}}d
       =1-i^{\otimes q}r^{\otimes q},\qquad
 r^{\otimes q}i^{\otimes q}=1,
\]
and $(h_q^{\mathrm{aux}})^2=r^{\otimes q}h_q^{\mathrm{aux}}
=h_q^{\mathrm{aux}}i^{\otimes q}=0$.
They commute with translation in each factor and with graded
permutations. Moreover $rj_s=1$ for every slot.
\end{lemma}

\begin{proof}
On generators, $Qs+sQ=N$ and $s^2=0$. The graded Leibniz
rule extends both identities to all PBW polynomials. Their
graded sums give the same identities on tensor powers, since
cross terms cancel. Translation preserves each auxiliary pair
and its number. Division by the positive integer $N_q$ gives
the contraction; on number zero, $i^{\otimes q}r^{\otimes q}$
is the identity. A nonzero value of $h_q^{\mathrm{aux}}$ has
positive auxiliary number, proving the side conditions.
The sums are invariant under graded permutations. Finally all
$U_s$ project to $J$ and all $-\eta_s$ to $c$.
\end{proof}

Write $h=h_1^{\mathrm{aux}}$. The maps $r,h$ are polynomial
state maps, not vertex morphisms. For example, auxiliary
fields killed by $r$ can have a singular product equal to the
vacuum. The collision-coefficient correction will account for
that product.

\subsection{Adjacent paths and all six pair lifts}

Let $A_r\subset\mathcal B$, for $r=1,2$, be generated by
\[
 U_r,\quad D_r=t_{r+1}-t_r,\quad
 p_r,p_{r+1},\eta_r,\eta_{r+1},\iota_r,\iota_{r+1}.
\]
Put $A=\operatorname{Com}(J+p_1+p_2,V_\kappa\otimes E^{\otimes2})$.
Each $A_r$ is a copy of this two-quartet commutant. The first
commutes with $j_3(\mathcal Z)$: it commutes with
$J+p_1+p_2$ and with the third quartet. The second is obtained
by sending $J$ to $J+p_1$ and shifting the two quartet indices;
it commutes with $j_1(\mathcal Z)$. The substitutions preserve
the generating brackets, the differential and the PBW basis.
Thus multiplication gives vertex maps
$A_1\otimes\mathcal Z\to\mathcal B$ and
$\mathcal Z\otimes A_2\to\mathcal B$.

Put polynomial forms $k[s,ds]$ before fields. In $A_r$ define
two paths on the generating states by
\begin{equation}\label{eq:vol2-supported-paths}
 \begin{aligned}
 F_r(J)&=U_r-s\kappa TD_r,&
 F_r(c)&=-\eta_r-s(\eta_{r+1}-\eta_r)-ds\,D_r,\\
 G_r(J)&=U_r-\kappa TD_r+sp_r+ds\,\iota_r,&
 G_r(c)&=-\eta_{r+1}.
 \end{aligned}
\end{equation}
Extend them by the vertex operations. Their images have current
self-bracket $\kappa\nu$, zero current--ghost bracket, and
closed ghost. The tensor differential gives
$dF_r(J)=\kappa TF_r(c)$ and $dG_r(J)=\kappa TG_r(c)$;
the parameter derivatives cancel the differentials of the
$ds$ terms. Hence they are vertex chain maps. They commute
with the fixed spectator throughout, by the commutant
calculation above. Their concatenation runs from $j_r$ to
$j_{r+1}$.

Let $h_r$ be the sum of their integrated $ds$ coefficients.
Each polynomial input has a finite polynomial integrand, so
integration is defined over $k$. Stokes' formula gives
\begin{equation}\label{eq:vol2-supported-path-integral}
 Qh_r+h_rd=j_{r+1}-j_r,
 \qquad h_r(J)=\iota_r,\quad h_r(c)=-D_r,\quad h_r(1)=0.
\end{equation}
The maps commute with translation and satisfy $rh_1=rh_2=0$.
Indeed every form coefficient contains an auxiliary factor
which survives the polynomial state projection until it is
set to zero.

Let $q_I:\mathscr R_I\to X^3$ be the ordered atlas of
Section~\ref{sec:vol2-action-carrier}. Put
$N_C=\omega_X\otimes C$ with
$(fa\,dz)\partial_z=-(f'a+fTa)\,dz$, and let
$P=q_I^!(N_C\boxtimes N_C\boxtimes N_C)$ be the regular
singleton source. For a partition $\pi$, let $\mathcal A_\pi$
be the external product of $\mathcal B^r=\omega_X\otimes\mathcal B$
over its block roots, localized where distinct roots are
separated and pushed onto the corresponding diagonal. This
direct image retains every normal derivative. Order the blocks
by their largest labels and transport their right densities
with that order. The three-position coefficient is
\begin{equation}\label{eq:vol2-supported-three-complex}
 \begin{split}
 \mathcal D_3&=\mathcal A_{1|2|3}
       \oplus\bigoplus_{ab}\mathcal A_{ab|c}[-1]
       \oplus\mathcal A_{123}[-2],\\
 D(x,y,z)&=(dx,b_1x-dy,b_2y+dz).
 \end{split}
\end{equation}
Here $b_1,b_2$ are full vertex contact maps, including
normal transfers and the Koszul signs. On regular states,
$b_1=\varrho_{12}+\varrho_{13}-\varrho_{23}$, where
\[
 \begin{aligned}
 \varrho_{12}(X,Y,W)&=[X_\lambda Y]\boxtimes W,\\
 \varrho_{13}(X,Y,W)&=(-1)^{|X||Y|}Y\boxtimes[X_\lambda W],\\
 \varrho_{23}(X,Y,W)&=X\boxtimes[Y_\mu W].
 \end{aligned}
\]
The second contacts are $[M_{\lambda+\mu}W]$ on $12|3$,
$[Y_\mu M]$ on $13|2$, and $[X_\lambda M]$ on $23|1$.
The vertex Jacobi identity gives
\[
 b_2b_1(X,Y,W)=[[X_\lambda Y]_{\lambda+\mu}W]
 +(-1)^{|X||Y|}[Y_\mu[X_\lambda W]]-[X_\lambda[Y_\mu W]]=0.
\]
These contact formulas extend by the differential-module
relations and retain their Laurent tests.

For $D_{ab}=q_I^{-1}\Delta_{ab}$ and
$D_{123}=q_I^{-1}\Delta_{123}$, including every ordered branch,
the support complexes are
\begin{equation}\label{eq:vol2-supported-pair-complexes}
 M_{ab}=R\Gamma_{D_{ab}}\mathcal D_3
       =\mathcal A_{ab|c}[-1]\oplus\mathcal A_{123}[-2],
 \qquad T_3=R\Gamma_{D_{123}}\mathcal D_3
       =\mathcal A_{123}[-2].
\end{equation}
To verify this description, apply the localization triangle to
the finite filtration by number of blocks. The separated term
is a direct image from the complement of $D_{ab}$ and has zero
extraordinary restriction there. The other pair terms meet
$D_{ab}$ only at the triple diagonal, away from which they are
localized. Only the two displayed terms remain. Their
inclusion is the support counit; localization agrees on chart
overlaps, so the description holds on the atlas.

An order $\sigma$ assigns its labels to slots $1,2,3$ and gives
a singleton insertion $E_\sigma:P\to\mathcal D_3$ through
$j_1,j_2,j_3$. The commuting slot algebras make its pair and
triple components zero. For a binary pair set
\[
 H_r(a,b)=-h_r(a)\boxtimes j_{r+1}(b)
                 +(-1)^{|a|}j_{r+1}(a)\boxtimes h_r(b).
\]
Its Hom differential is
$j_r\boxtimes j_{r+1}-j_{r+1}\boxtimes j_r$.
Use $H_1\boxtimes j_3$ for slots $1,2$ and
$j_1\boxtimes H_2$ for slots $2,3$, with the graded tensor
sign. Relabel source, target, coordinates and densities
together; call the resulting map $H_{\sigma,r}$.
The six edges are
\begin{equation}\label{eq:vol2-supported-hexagon-edges}
 \begin{array}{c|c|c|c}
 \text{edge}&\text{pair}&r&(\text{spectator},\text{slot})\\\hline
 123\to213&12&1&(3,3)\\
 213\to231&13&2&(2,1)\\
 231\to321&23&1&(1,3)\\
 321\to312&12&2&(3,1)\\
 312\to132&13&1&(2,3)\\
 132\to123&23&2&(1,1).
 \end{array}
\end{equation}
Both occurrences of $13$ exchange adjacent slots; neither
moves directly through a fixed middle slot.

\begin{proposition}\label{prop:vol2-supported-pair-lifts}
For an edge $e=(\sigma,r)$ exchanging labels $a,b$, put
\[
 \beta_e=-(b_1H_{\sigma,r})_{ab|c},\qquad
 \widehat\beta_e=(\beta_e,0):P\longrightarrow M_{ab}.
\]
This is a cochain map. Its composite with the support counit
is homotopic to $E_\sigma-E_{\sigma s_r}$ by
$(H_{\sigma,r},0,0)$. It retains the spectator and every normal
derivative. It vanishes on a vacuum in a paired input and
recovers the binary lift on a vacuum spectator.
\end{proposition}

\begin{proof}
The path equation gives the separated component of the
endpoint difference. The image of the paired homotopy lies
in $A_r$, which commutes with the spectator. Thus the two
unwanted first contacts vanish. The remaining contact is
\begin{equation}\label{eq:vol2-supported-binary-lift}
 \beta_r(a,b)=[h_r(a)_\nu j_{r+1}(b)]
       -(-1)^{|a|}[j_{r+1}(a)_\nu h_r(b)],
\end{equation}
tensored with the spectator and the prescribed sign. Its
image lies in $A_r$, so its second contact with that spectator
also vanishes. In the total complex this says
\[
 (E_\sigma-E_{\sigma s_r})
          -\delta(H_{\sigma,r},0,0)=(0,\beta_e,0).
\]
The right side is closed and factors through $M_{ab}$.

All maps commute with translation and use full contacts, so
their differential-module extensions retain the normal
derivatives. Vertex locality bounds the normal order on each
polynomial input. Relabelling transports the normal coordinate
and density once; the same formulas therefore agree on chart
overlaps and separation restrictions. Finally $h_r(1)=0$ and
the regular singular bracket with $1$ is zero, giving the
vacuum assertions.
\end{proof}

\subsection{The two-colour coefficient and its endpoints}

For free dg vertex algebras $D,E$, put
\[
 V(D,E)=j_*j^!(D^r\boxtimes E^r),\qquad
 W(D,E)=\Delta_*(D\otimes E)^r,
\]
where $j$ is the complement of the diagonal in the two-root
space and $D^r=\omega_X\otimes D$. If $v$ is the first root
minus the second, the tensor-product contact is
\begin{equation}\label{eq:vol2-coloured-contact}
 c_{DE}\bigl(f(v)(a\boxtimes b)\bigr)
 =\operatorname{pp}_v\bigl(f(v)(e^{vT_D}a)\otimes b\bigr).
\end{equation}
Only finitely many Taylor terms meet any fixed principal part.
The diagonal term retains the whole tensor $D\otimes E$.
This gives the two-colour complex
\[
 \mathcal Q(D,E)=V(D,E)\oplus W(D,E)[-1],\qquad
 D(x,z)=(dx,c_{DE}x-dz).
\]
Translation makes the contact differential-linear, and the
internal differential commutes with it.

For the resolving algebra itself let
$W(\mathcal B)=\Delta_*\mathcal B^r$ and put
\[
 \mathcal F_{\mathcal B}
   =V(\mathcal B,\mathcal B)\oplus W(\mathcal B)[-1],
 \qquad D(x,z)=(dx,b_{\mathcal B}x-dz),
\]
with its full vertex contact. Denote the positive diagonal
inclusion by $\iota$.

There are two orientations: $R$ for $(A,\mathcal Z)$ and
$L$ for $(\mathcal Z,A)$. Write $\alpha:A\to A_1$ and
$\gamma:A\to A_2$ for the two embeddings. Their placements
in the common resolution are
\begin{equation}\label{eq:vol2-coloured-placements}
 \begin{array}{c|cc}
 &u_+|_V&u_-|_V\\\hline
 R&\alpha\boxtimes j_3&\gamma\boxtimes j_1\\
 L&j_1\boxtimes\gamma&j_3\boxtimes\alpha.
 \end{array}
\end{equation}
On the diagonal use the corresponding normally ordered
products. The images commute, so these are vertex morphisms
and hence coefficient chain maps to $\mathcal F_{\mathcal B}$.

To use them at a pair $ab$, apply right differential-module
direct image along $i_{ab}:X^{\{ab,c\}}\hookrightarrow X^3$.
It includes the entire inner normal transfer and its density.
The subsequent diagonal direct image is $i_{123,+}$.
After the pair-support suspension, the placements have target
$M_{ab}$, and $\iota$ becomes the counit $T_3\to M_{ab}$.

Let $r_A:A\to C$ be the two-quartet polynomial retraction.
Projection in the $A$ factor gives maps $P_R=r_A\otimes1$
and $P_L=1\otimes r_A$, with components $P_V,P_W$.
Since $r_A$ commutes with translation,
\begin{equation}\label{eq:vol2-coloured-projection}
 P_Wc_\eta=c_{\mathcal Z\mathcal Z}P_V,
 \qquad \eta=R,L.
\end{equation}
This uses the tensor-product contact, not multiplicativity
of $r_A$ as a vertex map. The first-slot inclusion and the
binary auxiliary homotopy give a contraction $I,P,H_A$ of
these coloured complexes. The homotopy on the diagonal has
the negative of its state sign; equation
\eqref{eq:vol2-coloured-contact} makes the off-diagonal terms
cancel, so $\delta H_A=1-IP$.

On $V(\mathcal Z,\mathcal Z)$ define
\[
 H_1=-h_1\boxtimes j_2+j_2\boxtimes h_1,
 \qquad K_R=H_1+j_1\boxtimes h_2,
 \qquad K_L=H_1-h_2\boxtimes j_1.
\]
These tensor formulas include the graded signs. Extend them
over every finite meromorphic coefficient and make them zero
on the diagonal source. Their internal differentials are
$j_1\boxtimes j_3-j_2\boxtimes j_1$ and
$j_1\boxtimes j_2-j_3\boxtimes j_1$, respectively.
With $f_{st}(a,b)={:}j_sa\,j_tb{:}$, put
\begin{equation}\label{eq:vol2-coloured-defect}
 F_R=f_{13}-f_{21},\qquad F_L=f_{12}-f_{31},\qquad
 b_\eta=-b_{\mathcal B}K_\eta.
\end{equation}
Here $b_\eta$ has internal degree $-1$, while $F_\eta$ has
degree zero on the diagonal coefficient.

\begin{lemma}\label{lem:vol2-coloured-lift}
The full supported map is
\[
 \mathbb B^0_\eta(x,z)=b_\eta P_Vx+F_\eta P_Wz:
       \mathcal Q_\eta\longrightarrow W(\mathcal B)[-1].
\]
It is closed. If $E_A=u_+-u_-$, then
\begin{equation}\label{eq:vol2-coloured-endpoint-homotopy}
 E_A-\iota\mathbb B^0_\eta=\delta\mathbb K_A,
 \qquad \mathbb K_A=K_\eta P+E_AH_A.
\end{equation}
\end{lemma}

\begin{proof}
The path equation and the chain equations of the placements
give, in unshifted state degrees,
\begin{equation}\label{eq:vol2-coloured-chain-defect}
 \delta b_\eta=-F_\eta c_{\mathcal Z\mathcal Z},
 \qquad \delta F_\eta=0.
\end{equation}
Together with \eqref{eq:vol2-coloured-projection}, these are
the two components of closure into $W(\mathcal B)[-1]$.
On the two-singleton source, subtracting $\delta K_\eta$
from the endpoint difference leaves exactly its supported
components $b_\eta x+F_\eta z$. On arbitrary auxiliary
states, use $E_AI$ for that singleton difference and
$1-IP=\delta H_A$. This gives
\eqref{eq:vol2-coloured-endpoint-homotopy}, including the
term $E_AH_A$.
\end{proof}

The state projection has two useful identities:
\begin{equation}\label{eq:vol2-coloured-endpoint-projection}
 rF_R=rF_L=0,
 \qquad (r\boxtimes r)u_+|_V=(r\boxtimes r)u_-|_V=P_V.
\end{equation}
Indeed each placement is linear on free generators, and the
normally ordered product of its two commuting images is their
concatenated PBW polynomial. Both endpoints have the same
polynomial projection. This does not identify that polynomial
product with the vertex $(-1)$-product in $C$.
The second equality follows from $r\alpha=r\gamma=r_A$
and $rj_s=1$. Translation extends both to the complete
coefficient modules.

Define
\begin{equation}\label{eq:vol2-corrected-punctured-operation}
 s_\eta=hF_\eta,\qquad
 N_\eta=b_\eta+s_\eta c_{\mathcal Z\mathcal Z}.
\end{equation}
The auxiliary contraction and
\eqref{eq:vol2-coloured-endpoint-projection} give
$\delta s_\eta=F_\eta$; thus
\eqref{eq:vol2-coloured-chain-defect} gives $\delta N_\eta=0$.
On regular singleton inputs the tensor-product contact
vanishes. The additional $h_2$ term in $K_\eta$ has zero
regular contact with $j_1$, by the commutant property. Hence
$N_R=N_L=\beta_1$ there. On pole coefficients those terms
and the correction remain part of the operation.

\subsection{The outer map on the actual support}

Use the suspended coloured complex
\begin{equation}\label{eq:vol2-suspended-colours}
 \widetilde{\mathcal Q}_\eta
   =V_\eta[-1]\oplus W_\eta[-2],\qquad
 D(y,z)=(-dy,c_\eta y+dz).
\end{equation}
It is $\mathcal Q_\eta[-1]$ with its diagonal coordinate
multiplied by $-1$. Under the positive support counit, the
coloured lift becomes
\[
 \mathbb B_\eta(y,z)=-b_\eta P_Vy+F_\eta P_Wz.
\]

\begin{theorem}[Outer extension]
\label{thm:vol2-supported-outer-extension}
On the complete pair support define
\begin{equation}\label{eq:vol2-supported-outer-map}
 \mathcal C_\eta(y,z)=-N_\eta(r\boxtimes r)y,
 \qquad L_\eta(y,z)=-s_\eta P_Wz.
\end{equation}
The first is a degree-zero map $M_{ab}\to T_3$ and the second
has total degree $-1$ on the coloured source. Both placements
satisfy
\begin{equation}\label{eq:vol2-supported-placement-equation}
 \delta\mathcal C_\eta=0,
 \qquad \mathcal C_\eta u_\pm-\mathbb B_\eta=\delta L_\eta.
\end{equation}
Use $R$ for $12|3$ and $L$ for $23|1,13|2$, with the stated
root orders and normal direct images. If
$\mathbb K_A^{\mathrm{supp}}$ is the suspended homotopy of
\eqref{eq:vol2-coloured-endpoint-homotopy}, then
\[
 E_A-\iota\mathcal C_\eta u_\pm
 =\delta(\mathbb K_A^{\mathrm{supp}}-\iota L_\eta).
\]
Thus the extension retains the prescribed endpoint products.
\end{theorem}

\begin{proof}
The internal degree of $N_\eta$ is $-1$; the difference
between the pair and triple support shifts makes
$\mathcal C_\eta$ degree zero. It kills the diagonal component
of $M_{ab}$, so its chain equation is precisely
$\delta N_\eta=0$. This proves closure on arbitrary
polynomial auxiliary states and scalar poles.

On either placement, equations
\eqref{eq:vol2-coloured-projection} and
\eqref{eq:vol2-coloured-endpoint-projection} give the difference
\[
 -s_\eta P_Wc_\eta y-F_\eta P_Wz.
\]
The signs $(-,+)$ in
\eqref{eq:vol2-suspended-colours} and the odd degree of
$L_\eta$ give
\[
 \delta L_\eta(y,z)
 =-s_\eta P_Wc_\eta y-(Qs_\eta+s_\eta d)P_Wz,
\]
which is that difference. The final identity follows by
subtracting $\iota\delta L_\eta$ from the suspended endpoint
identity. Every equality is on the complete coloured source,
including its diagonal term.
\end{proof}

Neither placement needs to be inverted. In fact it is not an
equivalence. On diagonal state cohomology its map is
\[
 \Lambda(\alpha,\beta)\longrightarrow\Lambda(\varepsilon),
 \qquad \alpha,\beta\longmapsto\varepsilon.
\]
To compute these groups, contract the pairs
$(u_n,\kappa c_{n+1})$ in \eqref{eq:vol2-supported-state-complex};
only $c_0$ remains. The auxiliary contraction gives the same
cohomology for each resolving algebra. The displayed diagonal
map has kernel containing $\alpha-\beta$ and $\alpha\beta$.
Its separated part is a quasi-isomorphism, so the triangle of
the two-term coefficients identifies its cone with the shifted
cone of this diagonal map. Thus an inverse-placement method
would fail, while \eqref{eq:vol2-supported-outer-map} supplies
the outer extension directly.

The pole correction has a small explicit diagnostic. In the
$R$ orientation put
\[
 a=\iota_1+\iota_2,\quad p=p_1+p_2,\quad
 t=T(t_3-t_1),\quad e=T(\eta_3-\eta_1).
\]
For $x_J=v^{-1}(1\boxtimes J)$ and
$x_{Tc}=v^{-1}(1\boxtimes Tc)$, let $z_J,z_{Tc}$ denote the
corresponding diagonal transfers. The defining formulas give
\[
 \begin{gathered}
 Dx_J=\kappa x_{Tc}+z_J,\quad Dx_{Tc}=z_{Tc},\quad
 Dz_J=-\kappa z_{Tc},\\
 b_Rx_J=-a,\quad b_Rx_{Tc}=t,\quad
 F_Rz_J=p-\kappa t,\quad F_Rz_{Tc}=-e.
 \end{gathered}
\]
Since $h(p-\kappa t)=a$ and $h(-e)=-t$, the corrected
values are $N_Rx_J=-a+a=0$ and $N_Rx_{Tc}=t-t=0$.
Omitting the correction leaves chain defect $-(p-\kappa t)$
on $x_J$; reversing its sign leaves $-2(p-\kappa t)$.
These are nonzero states. They explain the correction on
poles without asserting an obstruction to the corrected map.

\section{The triple homotopy and its coefficient transfer}
\label{sec:vol2-supported-triple-homotopy}

The six pair lifts have two occurrences of each labelled pair.
A cube filling compares these occurrences. Its triple contact
currents sum to zero, whereas the actual iterated outer maps
need a further comparison. We construct both maps before
identifying their difference as a boundary.

\subsection{The nine-square filling}

Let $\mathsf I$ be the cellular chain complex of the interval
$[1,3]$, with vertices $v_1,v_2,v_3$ in degree zero, edges
$e_1,e_2$ in degree one and
$\partial e_r=v_{r+1}-v_r$. A cell of $\mathsf I^{\otimes3}$
assigns a height or a moving edge to each label, in label
order. Define
\[
 \Gamma(c_1\otimes c_2\otimes c_3)
 =\phi(c_1)\boxtimes\phi(c_2)\boxtimes\phi(c_3),
 \qquad \phi(v_s)=j_s,\quad\phi(e_r)=h_r.
\]
Its target is the separated term of
\eqref{eq:vol2-supported-three-complex}; a cell of dimension
$d$ gives an operation of degree $-d$. The graded tensor
formula and \eqref{eq:vol2-supported-path-integral} give
$\delta\Gamma(c)=\Gamma(\partial c)$ for the internal
differential.

For an edge $\sigma\to\sigma s_r$, let $x,y$ be its two
exchanged labels and $z$ the spectator at height $q$.
The corresponding chain is
\[
 c_{\sigma,r}=-(x:e_r,y:v_{r+1},z:v_q)
                  +(x:v_{r+1},y:e_r,z:v_q),
\]
with factors put in label order. Its boundary is the first
permutation vertex minus the second, and
$\Gamma(c_{\sigma,r})=H_{\sigma,r}$.
The sum around the hexagon is the boundary of
\begin{equation}\label{eq:vol2-nine-square-filling}
 \Sigma=\sum_{i=1}^3(-1)^i
       \bigl((e_1,e_2)_i+(e_2,e_1)_i+(e_2,e_2)_i\bigr),
\end{equation}
where $(c,c')_i$ fixes label $i$ at $v_3$ and assigns $c,c'$
to the other two labels in increasing order. To verify this,
the hexagon chain has twelve edges: it passes through the six
permutation vertices and through
\[
 (2,2,3),(3,1,3),(3,2,2),(3,3,1),(2,3,2),(1,3,3).
\]
On each of the three faces through $(3,3,3)$, it removes the
square at the opposite $(1,1)$ corner. The remaining three
squares are those in \eqref{eq:vol2-nine-square-filling}.
Their twelve interior edges cancel: two within each face and
two along each of the three shared edges. The exterior twelve
edges have precisely the hexagon orientation.

Put $\Phi=\Gamma(\Sigma)$ and define its contact maps
\begin{equation}\label{eq:vol2-cube-contacts}
 K_{ab}=(b_1\Phi)_{ab|c},\qquad \gamma_{ab}=-b_2K_{ab}.
\end{equation}
Their internal state degree is $-2$. Thus $(0,K_{ab},0)$
has total degree $-1$ and $(0,0,\gamma_{ab})$ has total
degree zero in \eqref{eq:vol2-supported-three-complex}.

\begin{proposition}\label{prop:vol2-cube-contacts}
If $e_+,e_-$ are the two hexagon occurrences of the pair $ab$,
then
\begin{equation}\label{eq:vol2-cube-pair-identity}
 \widehat\beta_{e_+}+\widehat\beta_{e_-}
       =\delta(0,K_{ab},0)+(0,0,\gamma_{ab}),
 \qquad \delta\gamma_{ab}=0,
\end{equation}
and $\gamma_{12}+\gamma_{23}+\gamma_{13}=0$.
The whole supported hexagon bounds:
\[
 \sum_e(H_e,0,0)+\sum_{ab}(0,K_{ab},0)=\delta(\Phi,0,0).
\]
\end{proposition}

\begin{proof}
The cube boundary gives
$\delta(\Phi,0,0)=(\sum_eH_e,b_1\Phi,0)$.
Differentiate it. The endpoint differences telescope, and
Proposition~\ref{prop:vol2-supported-pair-lifts} identifies
the remaining terms as the six supported lifts. Comparing
each pair component gives
$\beta_{e_+}+\beta_{e_-}=-\delta K_{ab}$ internally.
Its triple component gives the first identity with
$\gamma_{ab}=-b_2K_{ab}$. Differentiating that identity shows
$\gamma_{ab}$ is closed. Finally
$\sum\gamma_{ab}=-b_2b_1\Phi=0$ by vertex Jacobi.
\end{proof}

For the comparison with iterated brackets, set $g=h_1+h_2$
and define the merged-pair transport
\[
 k(a,b)=(-1)^{|a|}
   \bigl([h_1a_\nu h_2b]+[h_2a_\nu h_1b]+[h_2a_\nu h_2b]\bigr).
\]
Its internal differential is $\delta k=\beta_1-\beta_2$.
This follows by differentiating the two homotopies; the
contacts of $h_1$ with $j_3$ and of $h_2$ with $j_1$
vanish by the commutant property. Taking the contacts of the
nine squares gives
\begin{align}\label{eq:vol2-cube-pair-formulas}
 K_{12}(a,b,c)&=-k(a,b)\boxtimes j_3c
          +(-1)^{|a|+|b|}\beta_2(a,b)\boxtimes gc,\notag\\
 K_{23}(a,b,c)&=(-1)^{|a|}ga\boxtimes\beta_2(b,c)
                                      +j_3a\boxtimes k(b,c),\\
 K_{13}(a,b,c)&=(-1)^{|b|(|a|+1)}gb\boxtimes\beta_2(a,c)
                         +(-1)^{|a||b|}j_3b\boxtimes k(a,c).
                         \notag
\end{align}
For the first formula, the face fixing the spectator at
$v_3$ contributes the three terms of $-k$. On the other two
faces the spectator moves along $g$; the $h_1$ contacts with
$j_3$ vanish and the rest give $\beta_2$. The same calculation
gives the second formula. Moving label $2$ in front of
label $1$ gives the third formula and its Koszul signs.

\subsection{The three actual compositions}

Define the bracket on states by
\[
 B_\nu(a,b)=(-1)^{|a|}r\beta_1(a,b),
 \qquad x(a,b)=r\beta_1(a,b).
\]
It has degree $-1$ and is the operation fixed by the two
polynomial paths and the PBW retraction. Its three channels
are
\begin{equation}\label{eq:vol2-supported-three-channels}
 \begin{split}
 L(a,b,c)&=B_{\lambda+\mu}(B_\lambda(a,b),c),\\
 R_B(a,b,c)&=B_\lambda(a,B_\mu(b,c)),\\
 N_B(a,b,c)&=B_\mu(b,B_\lambda(a,c)),
 \qquad \epsilon=(-1)^{(|a|+1)(|b|+1)}.
 \end{split}
\end{equation}
The Fourier variables belong to positions $1,2$ relative to
position $3$. Sesquilinearity sends their translations to
$T_1=-\lambda$, $T_2=-\mu$, $T_3=T+\lambda+\mu$.

Choose the edges
$e_{12}:123\to213$, $e_{23}:231\to321$ and
$e_{13}:213\to231$. Their placements are respectively
$u_{R,+}$, $u_{L,-}$ and $u_{L,+}$. Define literal compositions
through the pair supports by
\begin{equation}\label{eq:vol2-supported-iterated-maps}
 G_{12}=\mathcal C_R\widehat\beta_{e_{12}},\qquad
 G_{23}=-\mathcal C_L\widehat\beta_{e_{23}},\qquad
 G_{13}=\mathcal C_L\widehat\beta_{e_{13}}.
\end{equation}
Their source remains $P$. To compare with the coloured map,
let $a_e:P\to\widetilde{\mathcal Q}_\eta$ have the indicated
binary lift in $A$ tensored with the spectator, and zero
diagonal component. Its outer scalar coefficient is regular,
so $c_\eta a_e=0$, including after any inner normal transfer.
It is therefore closed, $u_ea_e=\widehat\beta_e$ and
$L_\eta a_e=0$. Equation
\eqref{eq:vol2-supported-placement-equation} becomes
\[
 \mathcal C_\eta\widehat\beta_e=\mathbb B_\eta a_e.
\]
This is equality of operations on the whole regular singleton
source.

Projection of the preferred pair components gives
\[
 \begin{array}{c|c|c}
 \text{pair}&(r\boxtimes r)\beta_e(a,b,c)&\text{outer variable}\\\hline
 12&x(a,b)\boxtimes c&\lambda+\mu\\
 23&-(-1)^{|a|}a\boxtimes x(b,c)&\lambda\\
 13&(-1)^{|a||b|+|b|}b\boxtimes x(a,c)&\mu.
 \end{array}
\]
The minus sign in the second row is its first-contact sign;
the third row includes moving the spectator to the first root.
Since $N_R=N_L=\beta_1$ on a regular outer coefficient,
\begin{align}\label{eq:vol2-supported-tree-formulas}
 G_{12}(a,b,c)&=-\beta_1(x(a,b),c),\notag\\
 G_{23}(a,b,c)&=-(-1)^{|a|}\beta_1(a,x(b,c)),\\
 G_{13}(a,b,c)&=-(-1)^{|a||b|+|b|}\beta_1(b,x(a,c)).\notag
\end{align}
The degree of $x(a,b)$ is $|a|+|b|-1$, so
\begin{equation}\label{eq:vol2-supported-channel-projections}
 rG_{12}=(-1)^{|b|}L,
 \qquad rG_{23}=-(-1)^{|b|}R_B,
 \qquad rG_{13}=(-1)^{|b|}\epsilon N_B.
\end{equation}
For example, the first sign is
$-(-1)^{|a|+|b|-1}(-1)^{|a|}=(-1)^{|b|}$;
the third uses
$-(-1)^{|a||b|+|a|}=(-1)^{|b|}\epsilon$.

\subsection{Comparing cube contacts with the bracket}

The comparison uses a polynomial projection $\pi$ retaining
$U_1,\eta_1,p_1,\iota_1$ and killing every other auxiliary
coordinate and its translates.

\begin{lemma}\label{lem:vol2-supported-pbw-reduction}
For $Z\in\mathcal B$ and $a,b,y\in C$,
\begin{gather*}
 r[Z_\nu j_1y]=r[\pi Z_\nu j_1y],\qquad
 r[j_1y_\nu Z]=r[j_1y_\nu\pi Z],\\
 \pi h=h\pi,\qquad \pi Y=\pi j_2(rY)\quad(Y\in A_2),
 \qquad \pi hj_2=\pi h_1,\\
 r[h_1a_\nu h_1b]=0,\\
 r\beta_1(a,b)=r[h_1a_\nu j_1b]
                   -(-1)^{|a|}r[j_1a_\nu h_1b].
\end{gather*}
\end{lemma}

\begin{proof}
Only $U_1,p_1$ contract with $U_1$, and only $\iota_1$
contracts with $\eta_1$. A coordinate killed by $\pi$ thus
survives every contraction with $j_1y$ and is killed by $r$.
This proves the first two equalities. The auxiliary reverse
derivation carries the killed coordinates among themselves,
and their absence leaves auxiliary number unchanged, proving
$\pi h=h\pi$. On the generators of $A_2$, both $\pi$ and
$\pi j_2r$ send $U_2$ to $U_1+p_1$, send the ghosts to
$\eta_1$, and kill the remaining auxiliary fields. Hence
they agree on all PBW polynomials.

For $\pi hj_2=\pi h_1$, the first path contributes a factor
of $a_2$, which $\pi$ kills. In the second path, a monomial
with $m$ current factors choosing $p_1$ has form coefficient
obtained by replacing one of them by $\iota_1$ and weighting
the other $m-1$ by $s^{m-1}$. Integration gives $1/m$,
exactly the auxiliary contraction of the same monomial in
$j_2$. The graded sign of the odd replacement is the same.

To prove $r[h_1a_\nu h_1b]=0$, split $h_1=h^F+h^G$ by
the two paths. Every monomial of $h^F$ contains one translate
of $D=t_2-t_1$; every monomial of $h^G$ contains one translate
of $\iota_1$. Two $F$ terms leave both $D$ factors uncontracted,
and two $G$ terms leave both $\iota_1$ factors. Their
projections vanish. For a mixed term, use the generating
polynomials represented coefficientwise by
\[
 a=\exp\Bigl(\sum_n\alpha_nu_n+\sum_n\theta_nc_n\Bigr),
 \quad
 b=\exp\Bigl(\sum_n\alpha'_nu_n+\sum_n\theta'_nc_n\Bigr),
\]
where $\alpha,\alpha'$ are even and $\theta,\theta'$ odd.
Only finitely many parameters are needed for any chosen
coefficient. If $A_s,B_t$ are the zero-form exponents of the
two paths, their form coefficients are
\[
 h^Fa=\int_0^1\Bigl(\sum_n\theta_nT^nD\Bigr)e^{A_s}ds,
 \quad
 h^Gb=\int_0^1\Bigl(\sum_m\alpha'_mT^m\iota_1\Bigr)e^{B_t}dt.
\]
For the projection to survive, $D$ must contract with a
$tp_1$ factor of $B_t$ and $\iota_1$ with an
$-(1-s)\eta_1$ factor of $A_s$. Since
$[D_\nu p_1]=[\eta_1{}_\nu\iota_1]=1$, both contractions
have jet kernel
\[
 P_{nm}(z-w)=(-1)^n(n+m)!(z-w)^{-n-m-1}.
\]
Their product is $-t(1-s)$ times the square of the odd
element $\sum_{n,m}\theta_n\alpha'_mP_{nm}$, hence zero.
All other contractions multiply this zero. Reversing the
two paths gives the same argument with $\theta'$.
Taking polynomial coefficients proves the identity on all
states.

Finally, the contact of $h_1\otimes h_1$ followed by $r$
is zero. Its differential, using
$\delta h_1=j_2-j_1$, is the difference between
$r\beta_1$ and the last displayed bracket expression.
That difference is therefore zero, proving the final formula.
\end{proof}

Put
\begin{equation}\label{eq:vol2-cube-transfer-homotopy}
 T_{ab}=-r b_2(h\text{ on the merged factor})K_{ab},
 \qquad \widehat H=-(T_{12}+T_{23}+T_{13}).
\end{equation}
These maps have internal degree $-3$. The operation on the
merged factor uses its graded tensor sign.

\begin{proposition}\label{prop:vol2-cube-state-comparison}
On all polynomial states,
\[
 r\gamma_{12}=(-1)^{|b|}L+\delta T_{12},\qquad
 r\gamma_{23}=-(-1)^{|b|}R_B+\delta T_{23},\qquad
 r\gamma_{13}=(-1)^{|b|}\epsilon N_B+\delta T_{13}.
\]
Thus $\delta\widehat H=(-1)^{|b|}(L-R_B+\epsilon N_B)$.
\end{proposition}

\begin{proof}
For $12$, use $\delta(h\otimes1)=(1-j_1r)\otimes1$
and the fact that $rb_2$ commutes with the internal
differential. Equation \eqref{eq:vol2-cube-pair-identity}
gives
\[
 r\gamma_{12}-\delta T_{12}
 =-rb_2(j_1r\otimes1)K_{12}
   +rb_2(h\otimes1)(\beta_{e_+}+\beta_{e_-}).
\]
Here $\beta_{e_+}+\beta_{e_-}=
\beta_1\otimes j_3-\beta_2\otimes j_1$, and
$r\beta_1=r\beta_2=x$, since the slot shift has the same
polynomial projection in both binary copies. In the first
term, the part with $j_3$ has zero contact with $j_1$;
in the part with $g$, only $h_1$ remains. In the second term,
$h\beta_1\in A_1$ commutes with $j_3$.
Lemma~\ref{lem:vol2-supported-pbw-reduction} replaces the
remaining contact of $h\beta_2$ with $j_1$ by that of
$h_1x$. The result is
\[
 -rb_2(j_1x\otimes h_1+h_1x\otimes j_1)
       =-r\beta_1\circ(x\otimes1)=(-1)^{|b|}L.
\]

For $23$, the tensor formulas are
$K_{23}=g\otimes\beta_2+j_3\otimes k$ and
$\beta_{e_+}+\beta_{e_-}=j_1\otimes\beta_2-j_3\otimes\beta_1$.
Apply $1\otimes h$ to the merged factor. The same commutant
vanishings and reduction give
\[
 r\gamma_{23}-\delta T_{23}
 =rb_2(j_1\otimes h_1x-h_1\otimes j_1x)
 =-r\beta_1\circ(1\otimes x)=-(-1)^{|b|}R_B.
\]
The tensor $1\otimes x$ contributes $(-1)^{|a|}$.
For $13$, transpose $a,b$ and $\lambda,\mu$ with input sign
$(-1)^{|a||b|}$. This is a chain map. The result is
$-(-1)^{|a||b|+|a|}N_B=(-1)^{|b|}\epsilon N_B$.
Summing the three identities and using $\sum\gamma_{ab}=0$
proves the assertion about $\widehat H$.
\end{proof}

\begin{theorem}[Supported triple homotopy]
\label{thm:vol2-supported-triple-homotopy}
The maps
\begin{equation}\label{eq:vol2-supported-comparison-primitive}
 U_{ab}=j_1T_{ab}+h(\gamma_{ab}-G_{ab})
\end{equation}
satisfy $\delta U_{ab}=\gamma_{ab}-G_{ab}$. Consequently
\begin{equation}\label{eq:vol2-supported-triple-primitive}
 \mathcal H=-\sum_{ab}U_{ab}
       =j_1\widehat H+h(G_{12}+G_{23}+G_{13})
\end{equation}
has total degree $-1$ and
\[
 \delta\mathcal H=G_{12}+G_{23}+G_{13},\qquad
 r\mathcal H=\widehat H.
\]
These are maps into the full triple-supported coefficient
on every ordered branch, with all normal derivatives retained.
\end{theorem}

\begin{proof}
Both $\gamma_{ab}$ and $G_{ab}$ are closed. Their projection
difference is $\delta T_{ab}$ by
\eqref{eq:vol2-supported-channel-projections} and
Proposition~\ref{prop:vol2-cube-state-comparison}. The
auxiliary contraction gives
\[
 \delta\bigl(h(\gamma_{ab}-G_{ab})\bigr)
 =\gamma_{ab}-G_{ab}-j_1\delta T_{ab}.
\]
Adding $\delta(j_1T_{ab})$ proves the first equation.
The sum of the cube contacts is zero, so summation gives
\eqref{eq:vol2-supported-triple-primitive} and its differential.
The side condition $rh=0$ gives its projection.
Every constituent commutes with translation and extends over
the full normal transfer; all are polynomial on each input.
Thus these are identities of supported differential-module
operations, rather than only identities of state values.
\end{proof}

Changing a preferred occurrence of a pair changes its channel
by a specified boundary. Applying the closed outer map to
\eqref{eq:vol2-cube-pair-identity}, whose triple component it
kills, gives
\[
 \mathcal C_{ab}\widehat\beta_{e_+}
 +\mathcal C_{ab}\widehat\beta_{e_-}
 =\delta\bigl(\mathcal C_{ab}(0,K_{ab},0)\bigr).
\]
Change $\mathcal H$ by the same signed map. Changes for
different pairs add and therefore commute. This retains all
eight choices of one occurrence per pair, including both
$13|2$ edges.

\subsection{Normal transfers, substitutions and vacua}

The diagram just constructed has objects $P$, the coloured
complexes $\widetilde{\mathcal Q}_\eta$, the three pair
supports $M_{ab}$, the triple support $T_3$ and $\mathcal D_3$.
Its arrows include the first lifts $a_e$, both placements,
the outer maps and the support counits. Equation
\eqref{eq:vol2-supported-placement-equation} supplies the
placement homotopy, and
Theorem~\ref{thm:vol2-supported-triple-homotopy} supplies its
three-channel comparison.

If a specified collision substitution gives a chain map
$v:D'\to\widetilde{\mathcal Q}_\eta$, precomposition gives
\[
 \mathcal C_\eta u_\pm v-\mathbb B_\eta v
          =\delta(L_\eta v).
\]
For a further substitution $w$, the homotopies agree
literally: $(L_\eta v)w=L_\eta(vw)$. At three positions the
maps in this diagram are the vertex placements and the first
lifts already constructed. This statement does not assume
an unspecified system of higher substitution maps.

Normal differentiation acts on coefficients as well as
states. With $\nu$ the first-root Fourier variable and $T$
output translation, the corrected operation satisfies, for
every integer pole order $m$,
\begin{align}\label{eq:vol2-outer-normal-relations}
 N_\eta(Ta,b;v^{-m})
   &=-\nu N_\eta(a,b;v^{-m})
                       +mN_\eta(a,b;v^{-m-1}),\notag\\
 N_\eta(a,Tb;v^{-m})
   &=(T+\nu)N_\eta(a,b;v^{-m})
                       -mN_\eta(a,b;v^{-m-1}).
\end{align}
Differentiate $v^{-m}Y(a,v)b$ and then take principal parts
to obtain these identities for the contact part. The same
calculation with \eqref{eq:vol2-coloured-contact} gives them
for the correction. The initial closed embedding retains
the inner collision, the two-term coloured differential
retains the outer one, and the triple comparison retains their
intersection. Open restrictions localize these identities.
Relabelling transports states, normals and densities together.
Different ordered branches of the triple diagonal have no
overlap, so the formulas glue branchwise.

On a pole coefficient with either colour equal to the vacuum,
$rb_\eta=0$: the contact is a principal part of a translate of
$h_s(a)$, and $rh_s=0$. Also $rs_\eta=0$. Thus $N_\eta$
is closed with zero projection on that face, and $hN_\eta$
is its explicit null-homotopy. The canonical outer map has
the negative of this homotopy. On the common all-vacuum face
both are zero. The endpoint-relative homotopy remains the
one of Theorem~\ref{thm:vol2-supported-outer-extension}.

For the ternary maps, a vacuum in the paired inputs makes
the first lift zero. A vacuum spectator makes the regular
outer bracket zero. Hence every $G_{ab}$ vanishes on a
vacuum input. The formulas for $K_{ab}$ show that each
$\gamma_{ab}$ and $T_{ab}$ also vanishes there, so $U_{ab}$
and $\mathcal H$ do as well. A vacuum spectator still leaves
$K_{12}(a,b,1)=-k(a,b)\boxtimes1$; its outer contact is zero,
while its lower binary transport is retained.

\subsection{Removing the auxiliary fields on coefficients}

The state retraction must be corrected before it acts on
collision coefficients. For a finite set $I$ of $n$ positions,
write $\mathcal D_{A,I}=C_{A,I}[-n]$, for $A=\mathcal B$
or $\mathcal Z$. A partition with $q$ blocks has term
$\mathcal A_\pi[q-n]$. Split its differential as
\[
 D_A=D_{0,A}+b_A,\qquad
 D_{0,A}|_\pi=(-1)^{n-q}d_q,
\]
where $b_A$ merges blocks by full contacts.
Extend $i^{\otimes q},r^{\otimes q},h_q^{\mathrm{aux}}$
over the scalar coefficients and normal direct images.
On that partition put
\[
 I_\pi=i^{\otimes q},\qquad P_\pi=r^{\otimes q},\qquad
 H_\pi=(-1)^{n-q}h_q^{\mathrm{aux}}.
\]
Then $D_{0,\mathcal B}H+HD_{0,\mathcal B}=1-IP$.
The inclusion is a vertex map, so $b_{\mathcal B}I=Ib_{\mathcal Z}$.
Define finite operators
\begin{equation}\label{eq:vol2-supported-coefficient-transfer}
 R_I=\sum_{j=0}^{n-1}(-b_{\mathcal B}H)^j,
 \qquad \Pi_I=PR_I,\qquad \mathcal H_I=HR_I.
\end{equation}
Every contact lowers the block number. Thus
$(b_{\mathcal B}H)^n=0$ and $R_I=(1+b_{\mathcal B}H)^{-1}$.

\begin{proposition}\label{prop:vol2-supported-coefficient-transfer}
The maps in \eqref{eq:vol2-supported-coefficient-transfer}
give a contraction of the actual coefficient complexes:
\[
 \Pi_ID_{\mathcal B}=D_{\mathcal Z}\Pi_I,\qquad
 \Pi_II=1,\qquad
 D_{\mathcal B}\mathcal H_I+\mathcal H_ID_{\mathcal B}=1-I\Pi_I,
\]
with $\mathcal H_I^2=\Pi_I\mathcal H_I=\mathcal H_II=0$.
They preserve every fixed collision support and its counit,
commute with relabelling and open restriction, and preserve
the scalar vacuum maps. Their values have finite normal order
on each polynomial input.
\end{proposition}

\begin{proof}
Split the internal complex as $I\mathcal D_{\mathcal Z,I}
\oplus K$, where $K=\ker P$. On $K$ let the internal
differential be $d_K$ and the homotopy be $h_K$. Then
$d_Kh_K+h_Kd_K=1$, $h_K^2=0$. Since $I$ preserves contacts,
the full differential has the block form
\[
 D_{\mathcal B}=\begin{pmatrix}D_{\mathcal Z}&a\\0&D_K\end{pmatrix},
 \qquad D_K=d_K+b_K,\qquad D_{\mathcal Z}a+aD_K=0.
\]
The contact identities give $b_K^2=0$ and
$d_Kb_K+b_Kd_K=0$.
Put $A_K=1+b_Kh_K$. Direct multiplication gives
$D_KA_K=A_Kd_K$, since their difference is
$b_K(1-d_Kh_K-h_Kd_K)=0$.
Therefore $h'_K=h_KA_K^{-1}$ satisfies
$D_Kh'_K+h'_KD_K=1$.
Also $A_K^{-1}h_K=h_K$, so $(h'_K)^2=0$.

The displayed finite formulas are exactly
\[
 \Pi_I(z,k)=z-ah'_Kk,\qquad
 \mathcal H_I(z,k)=(0,h'_Kk).
\]
The asserted equations follow by matrix multiplication,
using $D_{\mathcal Z}a+aD_K=0$. In particular the target
differential is the original one with its vertex contacts.

The state maps commute with translation in each block and
are symmetric under graded permutations. The contacts have
the same equivariance with their Chevalley signs. Thus every
term agrees on ordered chart overlaps and open restrictions.
A partition lies in a fixed collision support when it is
coarser than its defining partition. State maps preserve
partitions and contacts only coarsen them, proving support
compatibility. The localization description of
\eqref{eq:vol2-supported-pair-complexes} identifies its
inclusion with the counit. The vacuum map factors through
$I$, on which $H$ vanishes. Finally there are at most $n-1$
contacts, each of finite normal order on a polynomial input.
\end{proof}

At three positions, $\Pi_3=P-PbH+PbHbH$.
For $w=\iota_1\boxtimes\eta_1\boxtimes1$ and
$x=p_1\boxtimes\eta_1\boxtimes1$, the differential is
$D_{\mathcal B}w=x+v_{\mathcal B}$, where
$v_{\mathcal B}$ is the simple pair-contact vacuum.
One has $Hx=w$, $Hw=Hv_{\mathcal B}=0$, and
$Pw=Px=0$, $Pv_{\mathcal B}=v_{\mathcal Z}$.
Consequently $\Pi_3(x)=-v_{\mathcal Z}$ and
$\Pi_3(v_{\mathcal B})=v_{\mathcal Z}$.
The uncorrected projection has defect $v_{\mathcal Z}\ne0$
on $w$.

On a pair support the formula becomes
\begin{equation}\label{eq:vol2-supported-pair-transfer}
 \Pi_M(y,z)=\bigl((r\boxtimes r)y,
                       rz+rb_2h_2^{\mathrm{aux}}y\bigr).
\end{equation}
There is at most one contact, and the pair homotopy has sign
$-h_2^{\mathrm{aux}}$, giving the plus sign. On the triple
support, $\Pi_T=r$ and $I_T=j_1$.
Let $I_M$ be its pair inclusion and put
\[
 \beta_e^C=\Pi_M\widehat\beta_e,\qquad
 \mathcal C_{ab}^C=\Pi_T\mathcal C_{ab}I_M.
\]
Then
\begin{equation}\label{eq:vol2-supported-composition-transfer}
 \Pi_T\mathcal C_{ab}\widehat\beta_e
                  =\mathcal C_{ab}^C\beta_e^C,\qquad
 \Pi_T\mathcal H=\widehat H.
\end{equation}
Indeed $\mathcal C_{ab}$ kills the triple component of the
pair homotopy, and $(r\boxtimes r)h_2^{\mathrm{aux}}=0$.
Hence $\mathcal C_{ab}\mathcal H_M=0$. Apply the closed
$\mathcal C_{ab}$ to $1-I_M\Pi_M=\delta\mathcal H_M$ to
obtain $\mathcal C_{ab}=\mathcal C_{ab}I_M\Pi_M$.
This proves the strict composition identity. The pair-to-triple
term of \eqref{eq:vol2-supported-pair-transfer} still occurs in
each transferred first lift and its endpoint homotopy.

\subsection{A nonzero state Jacobiator}

The homotopy is necessary even though the cube-contact sum
vanishes. Write $u=u_0$, $c=c_0$. The bracket fixed above
satisfies
\begin{align}\label{eq:vol2-composite-bracket-values}
 B_\nu(u^2,cTc)&=-2uTc+2\nu uc+2(Tu)c,\notag\\
 B_\nu(u^2,uc)&=-2u^2-\tfrac\kappa2\nu^2,\notag\\
 B_\nu(u^2,uTc)&=-2\nu u^2-2uTu-\tfrac\kappa6\nu^3,\\
 B_\nu(u^2,(Tu)c)&=-2uTu-\tfrac\kappa3\nu^3,\qquad
 B_\nu(u^2,u^2)=0.\notag
\end{align}
For the scalar terms, the parts of $h_1(uc)$, $h_1(uTc)$
and $h_1((Tu)c)$ without $\iota_1$ are respectively
\[
 -{:}U_1D_1{:}+\tfrac\kappa2{:}(TD_1)D_1{:},\quad
 -{:}U_1TD_1{:}+\tfrac\kappa2{:}(TD_1)^2{:},\quad
 -{:}(TU_1)D_1{:}+\tfrac\kappa2{:}(T^2D_1)D_1{:}.
\]
The factor $1/2$ is the integral of the first path parameter.
Every omitted term contains $\iota_1$, which cannot contract
with the second embedding and is killed by $r$.
The contractions of $U_2$ with $D_1,TD_1,T^2D_1$ are
$-v^{-1},-v^{-2},-2v^{-3}$. Their double contractions with
$U_2^2$, before the bracket's minus sign, are
$\kappa v^{-3},\kappa v^{-4},2\kappa v^{-4}$.
Converting $v^{-m-1}$ to $\nu^m/m!$ gives the three scalar
coefficients. Single contractions and the Taylor terms give
the displayed state terms. The two ordered ghost contractions
give the first line, and every term in the last bracket
retains an auxiliary field.

Substitution in the three channels gives
\[
 L(u^2,u^2,cTc)-R_B(u^2,u^2,cTc)
                   +\epsilon N_B(u^2,u^2,cTc)
                  =\frac\kappa3(\lambda+\mu)^3\mathbf1.
\]
The state terms cancel and $\epsilon=-1$.
Set $H(a,b,c)=(-1)^{|b|}\widehat H(a,b,c)$.
The ordinary Hom equation becomes
\begin{align}\label{eq:vol2-state-jacobi-homotopy}
 L-R_B+\epsilon N_B={}&dH(a,b,c)+H(da,b,c)\notag\\
 &-(-1)^{|a|}H(a,db,c)
             +(-1)^{|a|+|b|}H(a,b,dc).
\end{align}
Since $d(u^2)=2\kappa uTc$ and $d(cTc)=0$, degree and
the displayed value force
\[
 H(uTc,u^2,cTc)-H(u^2,uTc,cTc)=\frac{(\lambda+\mu)^3}{6}.
\]
The scalar vacuum need not be exact in $C$: the input
differential contributes the boundary in this operation
complex. The noncocycle example therefore does not define
an obstruction class from three cohomology classes.

The construction supplies outer operations on the complete
pair-support complexes and their triple homotopy for the
regular singleton source. Extending these fixed maps to
higher merged collision diagrams remains a further equation.
It must preserve the endpoint homotopies, proper collision
strata and vacuum faces already specified here.

\section{Directed composition beyond the associator}
\label{sec:vol2-directed-higher-composition}

A finite directed category isolates the next coherence
equation. Four arrows have five parenthesizations, so their
associativity homotopies must satisfy a pentagon relation.
This calculation fixes a carrier of rank one on every forward
arrow. Its solution will not identify that carrier with a
supported field-operation complex.

Fix objects $0,\ldots,4$, degree-zero arrows $a_{ij}$ for
$i<j$, scalar identities and no backward arrows. Products
are written chronologically. Introduce ten degree-zero
coefficients and their associativity equations:
\[
 a_{ij}a_{jk}=m_{ijk}a_{ik},\qquad
 F_{ijkl}=m_{ijk}m_{ikl}-m_{jkl}m_{ijl}.
\]
The dg coefficient algebra
\[
 K=k[m_{ijk}]\otimes\Lambda(\chi_{ijkl}),\qquad
 |\chi_{ijkl}|=-1,\quad d\chi_{ijkl}=F_{ijkl},\quad dm_{ijk}=0
\]
retains a chosen homotopy for each binary associativity
equation. All index sequences are increasing.

\begin{proposition}[The relative pentagon class]
\label{prop:vol2-directed-pentagon}
The element
\begin{align}\label{eq:vol2-directed-pentagon}
 S={}&m_{034}\chi_{0123}+m_{123}\chi_{0134}
           +m_{014}\chi_{1234}\notag\\
    &-m_{234}\chi_{0124}-m_{012}\chi_{0234}
\end{align}
is a nonzero class in $H^{-1}(K)$. Adjoining an even generator
$\zeta$ of degree $-2$ with $d\zeta=S$ therefore fills a
relation not already filled in $K$.
\end{proposition}

\begin{proof}
The five parenthesizations of the four-arrow path have
coefficients
\[
 \begin{array}{lll}
 P_1=m_{012}m_{023}m_{034},&P_2=m_{123}m_{013}m_{034},&
 P_3=m_{123}m_{134}m_{014},\\
 P_4=m_{234}m_{124}m_{014},&P_5=m_{012}m_{234}m_{024}.&
 \end{array}
\]
The differentials of the five displayed terms in $S$ are
$P_1-P_2,P_2-P_3,P_3-P_4,P_4-P_5,P_5-P_1$, so $dS=0$.
Every element of $K^{-2}$ is a polynomial combination of
products $\chi_I\chi_J$. Its differential has coefficients
of $m$-degree at least two. The nonzero coefficients of $S$
have $m$-degree one, so $S$ cannot be a boundary.
Finally $d^2\zeta=dS=0$; its even degree requires the
polynomial algebra $k[\zeta]$.
\end{proof}

The obstruction fixes the lower homotopies. In
$A=\Lambda(\varepsilon)$ with $|\varepsilon|=-1$ and $d=0$,
set every $m_{ijk}=1$, set $\chi_{0123}=\varepsilon$ and
all other $\chi$ to zero. All $F_{ijkl}$ vanish, giving a
map $K\to A$. It cannot extend to $K[\zeta]$, since $S$
maps to $\varepsilon$ and $A^{-2}=0$. Setting every $\chi$
to zero would extend. Thus the example obstructs prescribed
homotopies, not all coherent structures with those binary
products.

\subsection{The coefficient algebra in every arity}

For $N$ ordered objects $0,\ldots,N-1$, introduce a variable
$\beta_I$ of degree $2-n$ for every increasing sequence
$I=(i_0<\cdots<i_n)$, $n\geq2$. Let $R_N$ be their free
graded-commutative algebra, with
\begin{equation}\label{eq:vol2-directed-all-arity}
 d\beta_I=-\sum_{s=2}^{n-1}\sum_{r=0}^{n-s}
 (-1)^{r+(r+1)(2-s)}
 \beta_{(i_r,\ldots,i_{r+s})}
 \beta_{(i_0,\ldots,i_r,i_{r+s},\ldots,i_n)}.
\end{equation}
The inner coefficient is written first, fixing the sign
convention. The sum is empty on binary coefficients.
There are $2^N-1-N-\binom N2$ generators.

\begin{theorem}\label{thm:vol2-directed-all-arity}
Equation \eqref{eq:vol2-directed-all-arity} defines a
differential. For every commutative dg algebra $A$, maps
$R_N\to A$ specify strictly unital directed $A_\infty$
structures on the fixed labelled rank-one carrier, with
unary differential that of $A$. For $N=5$, the substitutions
\[
 \beta_{ijk}=m_{ijk},\qquad
 \beta_{ijkl}=-\chi_{ijkl},\qquad
 \beta_{01234}=-\zeta
\]
identify it with the preceding pentagon algebra.
\end{theorem}

\begin{proof}
Suspend each forward arrow to degree $-1$, denoting it by
$\bar a_{ij}$. Define the degree-one operations
\[
 b_n(\bar a_{i_0i_1}|\cdots|\bar a_{i_{n-1}i_n})
                  =\beta_I\bar a_{i_0i_n}.
\]
Keep coefficients on the left and extend them as a
coderivation of the tensor coalgebra over the object
idempotents. Its unary part differentiates coefficients.
Inserting an $s$-ary operation after $r$ inputs contributes
$(-1)^r$. Moving its coefficient of degree $2-s$ past those
inputs and then through the outer degree-one operation
contributes $(-1)^{(r+1)(2-s)}$. Thus the component of
$b^2$ with output length one is $d\beta_I$ plus the
insertion sum in \eqref{eq:vol2-directed-all-arity}.

Apply the prescribed derivation once more. Each resulting
term is a planar substitution with three operation vertices.
Its two internal edges can be contracted in exactly two
orders. For nested substitutions these are the two
associative orders of composing graded maps; for disjoint
substitutions they interchange two odd insertions. The
graded Leibniz rule gives opposite signs and the same
coefficient permutation in each pair. Hence every term
cancels and $d^2\beta_I=0$.

A map $R_N\to A$ therefore gives exactly the coefficients
with $b^2=0$. No nonidentity composable string has length
greater than $N-1$. Normalize away strings with an identity
input and impose the ordinary unit products; these give
the strictly unital structure. Conversely, its coefficients
give the map. Substitution at arities three and four gives
$d\chi=F$ and $d\zeta=S$ with the displayed signs.
\end{proof}

Polynomial differential forms on simplices can be used as
test algebras to retain paths and higher simplices of these
structures. The carrier and its bases stay fixed. Taking
a quotient by changes of basis or by general $A_\infty$
equivalences is a separate moduli problem. Nor does the
presentation assert that $R_N$ resolves the ordinary
associative-parameter ring at every singular point.

\subsection{An algebraic contraction of the invertible chart}

Assume $N\geq3$ and invert every binary coefficient:
\[
 R_N^\times=R_N[m_{ijk}^{-1}:i<j<k],\qquad
 q_{ij}=m_{0ij},\qquad
 L_N=k[q_{ij}^{\pm1}:1\leq i<j\leq N-1].
\]
These are algebraic localizations; no completion is used.
The first object supplies a contracting partner for every
coefficient on the other objects.

\begin{theorem}\label{thm:vol2-directed-invertible-chart}
There are dg algebra maps
$L_N\xrightarrow{\iota}R_N^\times\xrightarrow{p}L_N$,
with $p\iota=1$, and an $L_N$-linear degree-minus-one
homotopy satisfying $dH+Hd=1-\iota p$.
The inclusion sends $q_{ij}$ to $m_{0ij}$, and
\begin{equation}\label{eq:vol2-directed-retraction}
 p(m_{ijk})=\frac{q_{ij}q_{jk}}{q_{ik}}\quad(0<i<j<k),
 \qquad p(\beta_I)=0\quad(|I|\geq4).
\end{equation}
Thus $H^0(R_N^\times)=L_N$ and all other cohomology
vanishes, for every finite $N$.
\end{theorem}

\begin{proof}
For each increasing $I=(i_1,\ldots,i_\ell)$ with $i_1>0$
and $\ell\geq3$, set
\begin{equation}\label{eq:vol2-directed-contracting-pairs}
 \eta_I=q_{i_1i_\ell}^{-1}\beta_{0I},\qquad
 z_I=q_{i_1i_\ell}^{-1}d\beta_{0I}.
\end{equation}
Then $d\eta_I=z_I$ and $dz_I=0$.
In the insertion formula for $d\beta_{0I}$, the term
$r=1,s=\ell-1$ is $+\beta_Iq_{i_1i_\ell}$: the exponent
$1+2(3-\ell)$ is odd and the formula has an initial minus
sign. Every other coefficient with no zero index has fewer
than $\ell$ indices. Coefficients containing zero are
already $q$'s or unit multiples of $\eta$'s. Thus
$z_I=\beta_I+$ terms of smaller nonzero index size.
Increasing induction on $\ell$ inverts this triangular
polynomial change of generators.

For $I=(i,j,k)$,
$z_I=m_{ijk}-c_I$, with $c_I=q_{ij}q_{jk}/q_{ik}$ a
unit of $L_N$. Consequently $R_N^\times$ is the tensor
product over $L_N$ of the factors
\[
 A_I=
 \begin{cases}
 L_N[z_I,(z_I+c_I)^{-1}]\otimes\Lambda(\eta_I),&|I|=3,\\
 \operatorname{Sym}_{L_N}(L_Nz_I\oplus L_N\eta_I),&|I|\geq4,
 \end{cases}
\]
with $|z_I|=3-|I|$, $|\eta_I|=2-|I|$.
Evaluation at both variables zero is defined in every
factor, since each localized $c_I$ is a unit.

For $|I|=3$, a contraction of evaluation is
\[
 h(f(z))=\frac{f(z)-f(0)}z\eta,\qquad h(f(z)\eta)=0.
\]
The quotient belongs to the same localized ring because
the evaluation kernel is the ideal $(z)$.
For odd $|I|\geq5$, use
$h(z^a)=z^{a-1}\eta$ for $a\geq1$, $h(1)=0$ and
$h(z^a\eta)=0$. For even $|I|\geq4$, use
\[
 h(\eta^az)=\frac{\eta^{a+1}}{a+1},\qquad
 h(\eta^a)=0\quad(a\geq0).
\]
The graded Leibniz rule gives $dh+hd=1-\iota p$ in
each case, with $h\iota=ph=h^2=0$.
The denominators are invertible in characteristic zero.

Order the finite set of factors and take the signed
tensor homotopy
$H=\sum_j(\iota_1p_1)\otimes\cdots\otimes
(\iota_{j-1}p_{j-1})\otimes h_j\otimes1\otimes\cdots$.
The contraction equation telescopes, as in
Proposition~\ref{prop:vol2-independent-normal-support}.
After the binary substitutions, induction through the
triangular coordinates sends every higher coefficient
to zero. Evaluation is therefore exactly
\eqref{eq:vol2-directed-retraction}.
\end{proof}

The inclusion and projection are algebra maps; the
displayed homotopy is a homotopy of complexes. It also
determines the tangent complex before changes of arrow
basis. At the point with every $m_{ijk}=1$ and every higher
coefficient zero, let $T^\bullet$ be the dual of the
coefficient cotangent complex. Then
\begin{equation}\label{eq:vol2-directed-tangent-ranks}
 \dim T^j=\binom N{j+3},\qquad
 \operatorname{rank}(d:T^j\to T^{j+1})=\binom{N-1}{j+3}
 \quad(j\geq0),
\end{equation}
with out-of-range binomial coefficients zero. Indeed the
cotangent complex is obtained by differentiating the free
localized presentation. Each pair
\eqref{eq:vol2-directed-contracting-pairs} gives an identity
from tangent degree $|I|-3$ to $|I|-2$, and the
$\binom{N-1}{2}$ variables $q_{ij}$ remain unpaired.
Thus its cohomology is $k^{\binom{N-1}{2}}$ in degree zero
and zero otherwise. For $N=5$, the dimensions are
$(10,5,1)$ and ranks $(4,1)$; for $N=6$, they are
$(20,15,6,1)$ and $(10,5,1)$.

\subsection{Ordinary centres on two loci}

The tangent complex of the coefficient space differs from
the ordinary Hochschild complex of a category at a chosen
point. Let $C$ be the directed category on these $N$ objects,
with scalar identities and a single degree-zero basis arrow
$a_{ij}$ for each $i<j$.

\begin{proposition}\label{prop:vol2-directed-centre-loci}
At an associative point with every $m_{ijk}\ne0$,
$HH^0(C)=k$ and $HH^{>0}(C)=0$.
At the zero-product point,
\[
 \begin{gathered}
 \dim HH^0(C)=1,\qquad \dim HH^1(C)=\binom{N-1}{2},\\
 \dim HH^j(C)=\binom N{j+1}\quad(2\leq j\leq N-1),
 \end{gathered}
\]
and the higher groups vanish.
\end{proposition}

\begin{proof}
Associativity on $(0,i,j,k)$ gives
$m_{ijk}=q_{ij}q_{jk}/q_{ik}$. Replace $a_{ij}$ by
$q_{ij}^{-1}a_{ij}$ when $i>0$ and leave $a_{0i}$ unchanged.
Every composable product then has coefficient one.
The relative bar of Theorem~\ref{thm:vol2-directed-bar}
identifies its normalized Hochschild cochains with the
simplicial cochains of the ordered $(N-1)$-simplex.
For a positive-degree cochain put
\[
 (hf)(i_0,\ldots,i_{n-1})=
 \begin{cases}
 f(0,i_0,\ldots,i_{n-1}),&i_0>0,\\
 0,&i_0=0.
 \end{cases}
\]
In $dh+hd$, terms deleting any vertex other than the
inserted zero cancel. The remaining term is $f$ in positive
degree and $f-f(0)$ in degree zero. This contracts onto
the constants.

At zero products, every differential from degree one
upward is zero: each merger and endpoint action multiplies
two forward arrows. The degree-zero map is the difference
of vertex values on every forward edge. Its kernel is
$k$ and its rank $N-1$. Degree $j$ has one coordinate for
each increasing $(j+1)$-tuple. Subtracting that rank in
degree one gives the stated dimensions.
\end{proof}

At the zero-product point, the coefficient differential
has no linear part, so its tangent complex retains every
generator. Neither this tangent complex nor its ordinary
centre is obtained by substituting zero into the localized
contraction.

\section{The four-input comparison}
\label{sec:vol2-four-input-comparison}

The preceding constructions solve different composition
problems. The discarded coefficient complex retains outer
and crossed collision classes. The Heisenberg construction
gives actual pair-support outer maps and a triple homotopy
on regular singleton inputs. The directed coefficient
algebra retains all arities on its fixed rank-one carrier
and contracts on its invertible chart. A comparison with
field operations must supply the maps relating these
domains.

At four inputs, first construct a common operation complex
which retains the local vacuum and every nonconsecutive
collision channel. Transport the five parenthesizations
and their chosen associativity homotopies into it, including
the operations composing a homotopy with the remaining
inputs. Their pentagon then gives a specified cycle.
The unresolved equation is to construct a degree-minus-two
primitive of that cycle, or to prove that its class is
nonzero, with those lower homotopies fixed.

The directed algebra includes its higher generators and
their prescribed differentials. Theorem~\ref{thm:vol2-directed-invertible-chart}
computes its localized cohomology; it does not assert that
every choice of lower homotopies in another coefficient
algebra extends. The common supported operation complex and
the comparison into it remain to be constructed. Likewise
the three-channel scalar
current identity \eqref{eq:vol2-three-channel-current}
retains the contact terms needed by a field comparison,
but its residue relation alone is not that operation
homotopy. These missing maps preserve the distinction
between a coefficient calculation and a coherent action
on fields.
